\documentclass[12pt]{article}
\usepackage[utf8]{inputenc}
\usepackage[T2A,T1]{fontenc}
\usepackage{amssymb}
\usepackage{amsmath}
\usepackage{tempora}
\usepackage{csquotes}

\usepackage[english,russian]{babel}
\usepackage[backend=biber, style=ieee]{biblatex}

\textwidth 14cm
\textheight 21cm

\addbibresource{article_ru.bib}

\title{Несчётность множества континуум, шах}
\author{Роман Олегович Коротков}

\begin{document}

\selectlanguage{russian}

\begin{flushleft}
{\bf УДК 510.223 }\\{\bf MSC 03E99 }
\end{flushleft}

{\let\newpage\relax\maketitle}

\begin{abstract}

В статье критически пересматривается классическое доказательство несчётности множества вещественных чисел с помощью диагонального аргумента Кантора. Авторы возвращаются к наивной математике и разграничивают потенциальную и актуальную бесконечность, вводя два конструктивных множества на отрезке $[0, 1]$: потенциально бесконечное $C[0, 1]$ (объединение всех конечных двоичных записей) и актуально бесконечное $U[0, 1]$ (множество, в котором каждый разряд определён). Показано, что в случае потенциальной бесконечности диагональное построение либо не завершается, либо порождает лишь счётное множество пределов, не позволяя доказать существование числа вне исходного множества. В случае актуальной бесконечности применимость диагонального метода зависит от дополнительных предположений о структуре бесконечности - её универсальности, тотальности и степени рефлексивности. Различаются {\it сильная} и {\it слабая} интерпретации взаимно однозначного соответствия. Анализ оригинальной формулировки Кантора 1891 года показывает, что его рассуждение опирается на фиксацию универсальной вполне упорядоченной системы координат и не является общеприменимым вне этих предпосылок. Делается вывод, что диагональный аргумент сам по себе не даёт структурной информации о соотношении $\mathbb{N}$ и $\mathbb{R}$; его корректность определяется принятой морфологией бесконечности.

{\it Ключевые слова: диагональный метод Кантора; несчётность вещественных чисел; потенциальная и актуальная бесконечность; взаимно однозначное соответствие; сильная и слабая интерпретации взаимно однозначного соответствия; морфология бесконечности; наивная теория множеств; двоичные разложения; классы эквивалентности.}

\end{abstract}

\selectlanguage{english}

\begin{abstract}

The paper critically re-examines the classical proof of the uncountability of the continuum via Cantor’s diagonal argument. Returning to naive mathematics, the authors distinguish potential and actual infinity and introduce two constructive sets on the unit interval: the potentially infinite $C[0, 1]$ (the union of all finite binary expansions) and the actually infinite $U[0, 1]$ (the set in which every digit position is defined). It is shown that, under potential infinity, the diagonal construction either never terminates or yields only a countable collection of limits, and therefore cannot establish the existence of a real outside the original set. Under actual infinity the applicability of the diagonal method hinges on further assumptions about the morphology of infinity - its universality, totality and degree of reflexivity. A distinction is drawn between {\it strong} and {\it weak} interpretations of one-to-one correspondence. Analysis of Cantor’s original 1891 formulation reveals that the argument relies on the fixation of a universal, totally ordered coordinate system and is not generally valid outside these presuppositions. The authors conclude that the diagonal argument by itself supplies no structural information about the relation between $\mathbb{N}$ and $\mathbb{R}$; its correctness is determined by the adopted morphology of infinity.

{\it Keywords: Cantor’s diagonal argument; uncountability of the reals; potential and actual infinity; one-to-one correspondence; strong and weak interpretations of one-to-one correspondence; morphology of infinity; naive set theory; binary expansions; equivalence classes.}

\end{abstract}

\selectlanguage{russian}

Математические теории оперируют абстракциями и пока мы не выходим за их пределы все рассуждения совершенны, однако при попытке описать такими абстракциями реальные явления физического или воображаемого мира часто возникают трудности, более того порой даже попытки описать с их помощью объекты из другой области математики наталкиваются на нетривиальные препятствия. В частности применение современных теорий множеств, основанных на работах Георга Кантора, к арифметическим объектам требует переосмысления. Дело в том, что обе области неявно используют в своих рассуждениях бесконечные математические процедуры, под которыми мы понимаем неограниченно продолжаемые индуктивные или рекурсивные рассуждения. В свою очередь, сами бесконечные процедуры, как объекты математической теории, изучены недостаточно, что приводит к слишком вольному их применению. Счёт - есть ничто иное как бесконечная математическая процедура, неограниченный воспроизводимый эталон измерения количества элементов в коллекциях, которая связана с множеством натуральных чисел. Для записи одного единственного вещественного числа, а иногда и для его определения, также требуется такая процедура, взаимно однозначное соответствие, в свою очередь, описывается процедурой сопоставления элементов одного множества другому. Вопрос взаимодействия бесконечных математических процедур друг с другом нуждается в глубоком исследовании, особенно в контексте относительно новой области - когнитивной математики, в результате чего некоторые заключения могут оказаться уязвимыми.

Так, например, неясно как бесконечная процедура построения отличного от всех вещественного числа будет взаимодействовать с процедурами записи каждого из вещественных чисел и процедурой нумерации, по крайней мере, относительно такого взаимодействия можно сделать несколько различных предположений (что в равной степени применимо и к связи между производными математическими структурами). Выбор одного из таких предположений Кантором, без явного акцента и указания границ применимости, привёл, в конечном итоге, к неопределённости в доказательствах несчётности множества действительных чисел с помощью диагонального аргумента, изначально строгая формулировка после многократных воспроизведений утратила исходную чёткость. Нам предстоит начать с совершенно базовых понятий и, в рамках первых статей, оставить в стороне формализм теории множеств, то есть сделать наши рассуждения уязвимыми с точки зрения современной математики, такая уязвимость не означает некорректности, поскольку корректность математических рассуждений обусловлена связью с действительностью, а не соответствием некоторой системе принятых предположений. То есть пытаясь поставить формализм во главу угла, мы искажаем весь смысл и суть познания того мира, в котором мы существуем. Кроме того, вслед за Кантором нам придётся вернуться к рассмотрению концепции потенциальной и актуальной бесконечности, к которым в настоящее время крайне редко обращаются в математических статьях по теории множеств, несмотря на то, что как намереваются показать авторы, обоснование современной интерпретации теории множеств становится возможным только благодаря этим концепциям. Данные классические понятия окажутся не совсем удачными для описания механизмов возникновения интерпретаций теории множеств, тем не менее, прежде чем приступить к формированию нового языка, нам необходимо иметь хотя бы какое-то описание способное указать на причины возникновения неоднозначности. За что авторы просят прощения у дорожащих своим временем читателей, однако, иллюзорность некоторых предположений носит основополагающий характер и без тщательного разбора приводящих к ней построений может оставить сомнения в своей реальности, на что у авторов нет никакого права. Чтобы расставить все точки над i, нужно сказать ещё следующее, авторы не утверждают, что формализм теории множеств противоречив, однако, с уверенностью утверждают, что в нём присутствует неопределённость, связанная с трансцендентностью понятия бесконечности, результатом чего является наше непонимание этого формализма, именно оно приводит к неформально принятым соглашениям определяющим некоторую интерпретацию, которая, в свою очередь, рождает множественные курьёзы и парадоксы. Итак, формализм может быть уязвимым по двум причинам: содержать определённую неточность в отражении действительности сам по себе, или же наше понимание формализма может быть искажено связанными с ним методологиями и практиками. 

Необходимо сказать эти очевидные вещи, чтобы избежать необоснованной критики, так как авторы испытывают сомнения в том, что в рамках первой серии статей им удастся обрисовать альтернативное видение фундаментальных теоретико-множественных построений достаточно убедительно для того, чтобы помочь хотя бы части математического сообщества переключить своё восприятие с одной интерпретации на другую, поскольку такого рода попытки неоднократно предпринимались, наталкиваясь на враждебную позицию господствующих представлений основанных на вере в некоторые утверждения сопутствующие строгим аксиоматикам, но не являющиеся их частью. Впрочем, стоит также уточнить, что авторы не противопоставляют альтернативные подходы мейнстриму, предлагая общую рамку, где обе картины могут претендовать на корректность и сосуществовать дополняя друг друга. С  чего-то необходимо начинать, и именно взгляд под другим углом, долгое время подвергавшийся необоснованной цензуре \cite{transfinity}, противоречащей самой сути научного познания, заключающегося в беспрепятственном обмене опытом, позволяет рассмотреть тропинку к невероятно богатой и гармоничной математике, с позиций которой, в свою очередь, уже можно достоверно рассуждать о месте каждой из интерпретаций в структуре математических знаний. Поэтому авторы считают своим долгом зафиксировать результаты альтернативных подходов, содержащие все ключи к разрешению многих противоречий, продемонстрировав таким образом банальную истину: взаимопонимание и взаимоуважение оказываются зачастую продуктивней неприятия. Но прежде чем мы сможем подобраться к современной аксиоматической теории множеств и неформальным предположениям связанным с ней, необходимо вернуться к её истокам. Итак, в рамках первых двух статей мы отказываемся от формализма современной теории множеств и возвращаемся к классической наивной математике, но будем сопровождать наши рассуждения комментариями выражающими современный взгляд на описываемые проблемы. Рассмотрим сначала наиболее известное классическое доказательство несчётности множества вещественных чисел \cite{matan-diagonal}, основанное на опубликованной Кантором в 1891 году статье с доказательством теоремы о том, что множество объектов, зависимых от бесконечного количества координат, не может поставлено в соответствие множеству координат и следовательно имеет превосходящую мощность \cite{diagonal}, хотя исторически это и не первое доказательство несчётности предложенное им. Для простоты рассмотрения, а также для большего сходства с доказательством Кантора будем использовать двоичную систему исчисления для записи вещественных чисел.

Теорема: {\sl Множество $[0, 1]$ несчётно.}

Доказательство:

Каждое вещественное число $a \in [0, 1]$ имеет двоичную запись $a = 0.a_1a_2a_3\dots$, где $a_1$, $a_2$, $a_3$ – цифры $0$ или $1$. Здесь стоит сделать оговорку, одно вещественное число (в отличии от множества, рассматриваемого Кантором) в двоичной системе исчисления может быть записано двумя способами, например, $0.100(0)$ или $0.011(1)$. Так как равномощность множеству натуральных чисел подразумевает взаимно однозначное соответствие элементов, чтобы избежать неоднозначности будем использовать только первый вариант записи с $0$ для всех чисел, кроме $1 = 0.11(1)$. Предположим, что множество $[0, 1]$ счётно. Тогда его можно расположить в виде пронумерованной бесконечной последовательности:
\begin{align}
\begin{split}
a_1 &= 0.a^1_1 a^1_2 a^1_3 a^1_4\dots\\
a_2 &= 0.a^2_1 a^2_2 a^2_3 a^2_4\dots\\
a_3 &= 0.a^3_1 a^3_2 a^3_3 a^3_4\dots\\
\dots &= \dots\\
a_n &= 0.a^n_1 a^n_2 a^n_3 a^n_4\dots\\
\dots &= \dots
\end{split}
\end{align}

Построим новое число $b \in [0, 1]$ следующим образом: $b = 0.b_1b_2b_3\dots$, где $b_i$ – цифра отличная от $a^i_i$. Так как $b$ отличается от всех $a$ как минимум в одной цифре, то мы нашли число, которому не соответствует ни одно натуральное число, что противоречит исходному предположению. Теорема доказана.

Отметим, что мы избегаем употребления понятия биекция и говорим лишь о взаимно однозначном соответствии, таким образом, мы не отождествляем два этих понятия, биекция подразумевает, что образом одного множества является множество тождественное второму. Взаимно однозначное соответствие определяет правила поэлементного сопоставления, которые могут быть продолжены на все элементы обоих множеств, однако, продолженное до тождественности с одним множеством такое сопоставление не обязательно влечёт тождественность со вторым. Формально речь идёт о том, что функции определены и имеют смысл исключительно в рамках конкретной алгебры, в различных алгебраических системах функции имеющие одну и ту же форму могут существенно различаться как математические объекты. В классической онтологии понятие взаимно однозначного соответствия связано со структурными типами множеств - классами эквивалентности, а не с конкретными экземплярами множеств, таким образом, для установления взаимно однозначного соответствия достаточно воспроизводства структуры, а не полного отображения объектов друг на друга. Более того, само по себе взаимно однозначное соответствие в этой картине является классом эквивалентности множеств пар, а не фиксированным экземпляром такого множества. Будем называть интерпретацию отождествляющею биекцию и взаимно однозначное соответствие {\it сильной}, судя по всему, она является господствующей в рамках современной теории множеств, а альтернативную интерпретацию - {\it слабой}. Чтобы лучше понять разницу между ними, давайте рассмотрим совокупность натуральных чисел и совокупность чётных натуральных чисел делённых пополам. В рамках сильной интерпретации обе совокупности эквивалентны друг другу (тождественны в силу аксиомы экстенсиональности) и таким свойством - {\it рефлексивности} (бесконечности по Дедекинду) - мы наделяем все объекты, которые мы называем бесконечными множествами. В случае же слабой интерпретации две совокупности не тождественны друг другу, но обладают некоторой общностью - структурной неограниченностью, позволяющей нам проверить наличие сколь угодно больших элементов в обоих совокупностях, такое своеобразное отношение между совокупностями является отношением эквивалентности формирующим класс эквивалентности, который мы связываем с понятием мощности. Авторы считают, что Кантор в своих работах, использует именно слабую интерпретацию, которая на некотором историческом этапе, обращается в сильную во время синтеза аксиоматической теории множеств, что приводит к некоторым артефактам. Такое упрощение разделяет теорию Кантора и современную теорию, что существенно усложняет поиск уязвимостей полученного формализма. Ещё точнее, Кантор оперирует двумя видами сопоставления элементов: эквивалентностью или кардинальным соответствием и подобием или ординальным соответствием, они вводятся процедурно, как типовые операции над множествами, в то время как, в современной интерпретации оба понятия соотносятся с единой биекцией. В следующих статьях серии, мы обнаружим, что изоморфизм между современным подходом и подходом Кантора доказывается лишь для ординального соответствия, в то время как, взаимосвязь понятий мощности в этих подходах до сих пор сохраняет неопределённый статус. В результате проекция кардинальных рассуждений Кантора на современную теорию оказывается не до конца обоснованной. Это объясняется тем, что Кантор сформировался в рамках классической образовательной модели, где философии принадлежала интегрирующая роль, что и обусловило онтологическую и эпистемологическую согласованность его построений. В то время как фундамент современных теорий множеств закладывался в эпоху стремительной профессионализации и специализации математики, когда фокус сместился в сторону формальной строгости и прикладной эффективности. Данный сдвиг, маргинализировавший философскую рефлексию, привёл к систематической метафизической неопределённости и концептуальной незавершённости этих теорий. Итак, до принятия аксиоматик теории множеств, мы не можем однозначно интерпретировать взаимно однозначное соответствие тем или иным способом, поскольку не имеем на то достаточных оснований, в связи с чем мы будем исходить из того, что обе интерпретации допустимы. 

Рассмотрим некоторое множество вещественных чисел $U[0, 1]$ на отрезке $[0, 1]$, структурированное особым образом (вст. \ref{seс:binary-set}):
\begin{align}
\begin{split}
&0.000000000\dots\\
&0.100000000\dots\\
&0.010000000\dots\\
&0.110000000\dots\\
&\dots
\label{seс:binary-set}
\end{split}
\end{align}

Правила построения такого бесконечного множества очень просты. Начинаем с числа $0.0(0)$, объединяем его со множеством, получаем $U_0[0, 1]$, состоящее из $1$ элемента. Выбираем из множества действительных чисел на отрезке $[0, 1]$ все числа имеющие 1 в первом разряде (позиции), где ни одно из имеющихся на данный момент чисел не имеет $1$, или что тоже самое в разряде соответствующему номеру итерации, и имеющие $0$ во всех последующих позициях, получаем $0.10(0)$, объединяем новые числа со множеством, теперь $U_1[0, 1]$ состоит из двух элементов. На следующей итерации снова выбираем числа с $1$ в новом разряде, получаем $0.010(0)$ и $0.110(0)$, добавляем их ко множеству, получаем $4$ элемента в $U_2[0, 1]$. Тоже самое делаем на $i$-той итерации, получаем $U_i[0, 1]$ из $2^i$ элементов и т.д. У нас получилась некоторая бесконечная процедура (или трансфинитная, если угодно, нам достаточно продолжать её до первого трансфинитного порядкового числа $\omega$). Описанный процесс допускает трансфинитное продолжение, например, на итерации $\alpha \geq \omega$, подмножество $U_{\alpha}[0, 1]$ станет несчётным, тем не менее, мы можем выбрать соответствующий ему отрезок, сдвинуть на некоторую обобщённую величину $2^{-\alpha}$ и объединить оба отрезка. Данную процедуру можно определить и аналитически, поскольку любое вещественное число в бесконечной двоичной записи, по определению, есть сумма ряда \cite{algorithms}:
\begin{equation}
a = \sum_{i=1}^{\infty} a_i2^{-i}
\end{equation}

Здесь $a_i$ - коэффициенты $0$ или $1$ соответствующие $i$-тому разряду двоичной записи действительного числа. Тогда процесс построения можно переформулировать следующим образом. Начинаем, как и прежде с $a^0 = 0.00000\dots = 0$, $a^0_i = 0$ для всех $i$, имеем $U_0[0, 1]$. Добавляем ко всем элементам имеющимся на данный момент $2^{-1}$, получаем $a^1 = 0.10(0)$, $a^1_1 = 1$ и $a^1_i = 0$ для всех $i > 1$, объединяем новые ряды с $U_0[0, 1]$, имеем $U_1[0, 1]$ состоящее из двух рядов. И так далее, на $i$-той итерации добавляем слагаемое $2^{-i}$ в каждый из имеющихся в $U_{i-1}[0, 1]$ рядов, а полученные таким образом новые ряды, в свою очередь, объединяем со множеством, получаем $U_i[0, 1]$. Повторив такую процедуру бесконечное число раз мы получим бесконечное упорядоченное множество двоичных рядов $U[0, 1] = \bigcup\limits^{\infty}_{i = 0} U_i[0, 1] = U_\infty[0, 1]$. На каждой итерации количество элементов во множестве удваивается, значит при переходе $i$ к бесконечности, количество рядов равное $2^i$ также станет бесконечностью.

Можно заметить, что до перехода к бесконечности такая конструкция, являются зеркальным отражением чисел натурального ряда в двоичной записи, обозначим множество полученное таким способом как $C[0, 1]$, но поскольку в математической среде до настоящего времени сохранилась точка зрения о потенциальной бесконечности натурального ряда, отличающей его от множества вещественных чисел, расширение которого бесконечно удалёнными точками ничему не противоречит, то нам потребовались независимые от натурального ряда процедуры. Также для построения $C[0, 1]$ мы можем воспользоваться одним из указанных выше способов, но продолжать процесс не до бесконечности, а для неограниченного количества итераций $i \in \mathbb{N}$, то есть $C[0, 1] = \bigcup\limits_{i \in \mathbb{N}} C_i[0, 1]$, ведь натуральные числа в двоичной записи те же самые ряды только с положительными степенями двойки. Итак, у нас есть три способа задания похожих множеств, с помощью которых мы можем построить потенциально бесконечное множество $C[0, 1]$ и актуально бесконечное множество $U[0, 1]$. Необходимо отметить, что согласно современным представлениям, $U[0, 1] = \bigcup\limits^{\infty}_{i = 0} U_i[0, 1]$ является счётным множеством, поскольку представляет собой объединение всех вещественных чисел конечной длины, сколь угодно большой, но конечной (можно сказать, что это множество всех наследственно конечных последовательностей). Однако, это именно то, что обычно подразумевается под потенциальной бесконечностью, поэтому такое представление мы выделяем в объект $C[0, 1]$, который с современных позиций и должен определяться как $C[0, 1] = \bigcup\limits^{\infty}_{i = 0} C_i[0, 1]$. Множество $U[0, 1]$ отличается от $C[0, 1]$ тем, что оно актуально бесконечно, что проще всего представить себе так: на каждом шаге $i$ построения множества $U[0, 1]$ мы записываем цифру $0$ в $i$-ую позицию всех чисел, которые уже имеются во множестве, или добавляем в них слагаемое $0*2^{-i}$, а в копии - записываем цифру $1$ или добавляем $1*2^{-i}$. Можно также представить себе трансфинитную последовательность заготовок чисел все разряды которых $0$, в которую мы записываем поразрядно $1$, копируя начальные сегменты и спускаясь вниз по экспоненциальной лестнице, пока не пройдём все разряды. Таким образом, проделав такие операции для всех натуральных чисел, мы получаем множество $U[0, 1]$ состоящее из вещественных чисел, имеющих цифру или слагаемое в любой позиции соответствующей натуральному числу. Если представить данную конструкцию в матричном виде, то множество столбцов является множеством натуральных чисел, и все такие столбцы содержат поровну нулевых и единичных цифр или слагаемых. В самом деле, если предположить существование столбца $k$ с неопределённой цифрой или слагаемым, то мы проделали операцию не для всех натуральных чисел, а только для чисел $< k$. Согласно современной интерпретации теории множеств, несмотря на то, что множество столбцов имеет порядковый тип $\omega$ считается, что множество $U[0, 1]$ не содержит элементов с бесконечным определением, то есть возникает некоторый объект $U_{<\omega}[0, 1]$ содержащий все элементы с конечным количеством использованных разрядов (ведь $U_{\omega}[0, 1]$ по определению должен иметь позицию или слагаемое $\omega$, то есть множество столбцов должно соответствовать порядковому типу $\omega + 1$). Это может показаться странным, учитывая то, что мы определяем $\omega$ как $\bigcup\limits^{\infty}_{i = 0} i$, в частности, именно это мы наблюдаем в множестве столбцов. Ведь у нас просто нет какого-либо иного определения бесконечности - только через объединение всех конечных элементов. Тем не менее, данное представление является необходимым условием отсутствия противоречий в современной интерпретации. Поэтому эта странность объясняется тем, что не все бесконечные последовательности выразимы в теории и когда мы записываем бесконечное объединение, мы подразумеваем что оно состоит только из выразимых элементов. Здесь мы стремимся сделать минимум возможных предположений относительно природы бесконечности, поэтому рассматриваем два варианта, потенциально бесконечный $C[0, 1]$, с элементами со сколь угодно большой длиной, и актуально бесконечный $U[0, 1]$, с элементами бесконечной длины. Поскольку мы договорились работать в наивной математике нам не нужно, привлекать сложный аппарат нестандартного анализа или арифметики второго порядка для формализации определения $U[0, 1]$. С точки зрения современной теории $U[0, 1]$ можно ассоциировать с $2^{\aleph_0}$, а $C[0, 1]$ - с $2^{\omega}$, такая ассоциация весьма условна, поскольку, как мы выясним, некоторые предпринимаемые нами действия совместимы не со всеми интерпретациями ZFC и NBG, тем не менее, они совместимы с общепринятой.

К сожалению в математике отсутствует точное семантическое определение потенциальной и актуальной бесконечности (аксиома бесконечности определяет такой объект, но не наделяет его семантикой), поэтому мы будем исходить из того, что под потенциальной бесконечностью понимается свойство натурального ряда быть неограниченно продолжаемым, выраженное формально в принципе Архимеда: для всякого натурального числа существует большее натуральное число \cite{matan-archimed}, который так или иначе присутствует в любой модели натуральных чисел. То есть, мы не можем говорить о натуральном ряде как об объекте представляющем все натуральные числа, а подразумеваем под ним лишь совокупное название объектов получаемых определённым способом, что находит воплощение в множестве $C[0, 1]$. Что понимается под актуальной бесконечностью не всегда понятно и зачастую парадоксально, как, например, завершимость незавершимого, из чего вряд ли получится создать нечто непротиворечивое, но мы постараемся рассмотреть некоторые конструкции, которые связывают с этой загадочной сущностью. Формализм теории множеств принимает существование абсолютной бесконечности в качестве аксиомы, точнее он просто декларирует что есть какая-то бесконечность, без пояснения какая именно, но факт того что она является завершённым объектом теории, описывается именно термином актуальной бесконечности. Однако, следствия этого существования оказываются довольно неоднозначными. Поэтому для наших рассуждений, мы примем в качестве аксиомы менее формальное утверждение, что бесконечные математические процедуры (или бесконечные индуктивные построения) возможны, что воплощается во множестве $U[0, 1]$ и, как мы выясним, является необходимым для того, чтобы исследовать различные варианты применимости диагонального метода. Другой, более классический подход, заключается в использовании предельного перехода к бесконечности, что подразумевает концепцию сходимости к некоторому идеалу. Оба этих подхода могут быть применены для определения вещественных чисел. Так сам Кантор \cite{settheory} и Дедекинд \cite{irrationality} определяли вещественные числа используя предельные переходы, поскольку они пытались выразить их через рациональные числа, таким образом мы не идентифицируем само вещественное число, а как бы указываем на него с помощью последовательности, ведь вещественное число - предел последовательности не принадлежит самой последовательности, что легко видеть рассмотрев строго монотонную последовательность и предположив, что её предел имеет в ней какой-либо номер, поэтому сам предел характеризуют как актуальную бесконечность:
\begin{equation}
a = \lim_{n\to\infty} a_n, a \notin \{a_n\}
\end{equation}

Но вещественные числа можно определить и явно с помощью бесконечной записи по некоторому основанию, как сделал Вейерштрас \cite{decimal}, или бесконечного ряда:
\begin{align}
\begin{split}
1 = 0.111(1) &= \sum_{i=1}^{\infty} 2^{-i}\\
1/3 = 0.010101(01) &= \sum_{i=1}^{\infty} 2^{-2i-1}\\
\pi/4 = 0.110010010000111\dots &= \sum_{i=1}^{\infty} a_i^{\pi/4}2^{-i},\thickspace a_i^{\pi/4} \in \{0, 1\}
\end{split}
\end{align}

Так в первой строчке речь идёт именно о тождественной форме записи $1$ с помощью бесконечной процедуры, иначе, согласно апории Зенона движение было бы невозможно, поскольку чтобы преодолеть путь нужно сначала преодолеть половину пути, а чтобы преодолеть половину пути, нужно сначала преодолеть половину половины пути, и так далее, а Ахилес никогда бы не смог догнать черепаху, что может быть и можно вообразить, но наша вселенная, к счастью, умеет оперировать с бесконечностью. С формальной же точки зрения это простое следствие того, что мы определяем вещественные числа как классы эквивалентности объектов различие между которыми невыразимо никаким конечным (рациональным) числом. Следствие требования архимедовости для этих классов эквивалентности мы можем наблюдать, например, как наличие двух позиционных форм записи фиксированного числа. Обратим внимание на то, что здесь нам требуется предположение об алгебраической замкнутости бесконечных сумм, которое можно обосновать замкнутостью операции сложения, что выглядит более естественным, по сравнению с замыканием множества со всеми пределами внутренних последовательностей. Данное наблюдение возвращает нас к странности с интерпретацией $U[0, 1]$ в современном контексте, в самом деле, если в теории появляется бесконечность, она позволяет ввести бесконечные ряды и объединения, но тогда запись $\bigcup\limits^{\infty}_{i = 0} U_i[0, 1]$ семантически будет связана именно с бесконечностью в теории и независима от выразимости элементов в мета-теории, однако мы утверждаем существование в теории объектов, которые определимы только через ограничения мета-теории. Насколько такая интерпретация может быть обоснованной нам и предстоит выяснить. Когда мы используем бесконечную запись по какому-либо основанию или бесконечные ряды мы уже принимаем аксиому возможности бесконечных математических процедур (что можно интерпретировать как расширению множества натуральных чисел элементом $\infty$, $\mathbb{N}^{\infty} = \mathbb{N} + \infty$, но поскольку это противоречит господствующим представлениям теории множеств оставим пока этот вопрос в стороне). Для двоичной записи вещественных чисел нам требуются бесконечные математические процедуры, в свою очередь для построения диагонального элемента нам требуется бесконечная двоичная запись вещественных чисел. Выразимость в данном контексте нас не волнует, поскольку теория уже оперирует бесконечными суммами как объектами и гарантирует замкнутость, то есть факт того, что результат также является объектом теории даже в случае бесконечного количества изменений элементов суммы. Таким образом, то что мы воспользовались этими процедурами для построения $U[0, 1]$, произведя ровно столько итераций, сколько требуется для двоичной записи вещественных чисел, ни больше, ни меньше, ничему не противоречит. 

В построенных математической структурах $C[0, 1]$ и $U[0, 1]$ можно выделить некоторую последовательность элементов и оценить её предел при переходе количества итераций к бесконечности. Например, возьмём элемент $c_0 = u_0 = 0$, на следующей итерации - его же $c_1 = u_1 = 0.0$, а дальше возьмём элемент полученный из $0$ на второй итерации – $c_2 = u_2 = 0.01$, тоже самое проделаем со вторым элементом и получим третий - $c_3 = u_3 = 0.010$, далее четвёртый $c_4 = u_4 = 0.0101$. Повторяя такие действия неограниченное число раз в $C[0, 1]$ мы получим последовательность, предел которой равен $0.01(01)$, в свою очередь $U[0, 1]$ будет содержать и сам этот элемент:
\begin{equation}
\lim_{n \in \mathbb{N}} c^{0.01(01)}_n = 0.01(01) = u^{0.01(01)}_\infty = \sum_{i=1}^{\infty} 2^{-2i-1}
\end{equation}

Точно также мы можем выделить последовательность элементов значащие цифры которых совпадает с двоичной записью числа $\pi/4$, так на каждой итерации, если в двоичной записи числа $\pi/4$ стоит $0$ мы добавляем в последовательность предыдущий элемент, а если $1$ то добавляем элемент полученный на данной итерации из предыдущего: $c_0 = u_0 = 0$, $c_1 = u_1 = 0.1$, $c_2 = u_2 = 0.11$, $c_3 = u_3 = 0.110$, $c_4 = u_4 = 0.1100$, $c_5 = u_5 = 0.11001$, $\dots$. Продолжая неограниченно долго в $C[0, 1]$ мы получим последовательность, предел которой равен $\pi/4$, в $U[0, 1]$ бесконечный элемент последовательности будет самим числом $\pi/4$, поскольку по построению бесконечные двоичные записи чисел будут совпадать:
\begin{equation}
\lim_{n \in \mathbb{N}} c^{\pi/4}_n = \pi/4 = u^{\pi/4}_\infty = \sum_{i=1}^{\infty} a_i^{\pi/4}2^{-i},\thickspace a_i^{\pi/4} \in \{0, 1\}
\end{equation}

Будем обозначать неограниченные последовательности вещественных чисел, каждый следующий элемент которых имеет на один значащий разряд двоичной записи больше чем предыдущий, как $\{c_n\} \in C[0, 1]$ и $\{u_n\} \in U[0, 1]$.

Числа имеющие бесконечную двоичную запись, такие как $0.01(01)$ или $\pi/4$, в силу господствующих убеждений не могут быть поставлены во взаимно однозначное соответствие с элементами $C[0, 1]$. В $C[0, 1]$ содержатся только вещественные числа соответствующие натуральным, и уместно говорить лишь о том, что последовательности $\{c_n\}$ имеют некоторые пределы, сами эти пределы могут и не принадлежать множеству $C[0, 1]$. В $U[0, 1]$, в свою очередь, могут содержаться вещественные числа имеющие бесконечную двоичную запись. Действительно, если мы попытаемся вычислить порядковый номер каждого из этих чисел то получим:
\begin{align}
\begin{split}
n_{0.01(01)} &= \sum_{i=1}^{\infty} 2^{2i+1} = \infty\\
n_{\pi/4} &= \sum_{i=1}^{\infty} a_i^{\pi/4}2^{i} = \infty,\thickspace a_i^{\pi/4} \in \{0, 1\}
\end{split}
\end{align}

Рассмотрим теперь как будет взаимодействовать диагональный метод с построенными подмножествами вещественных чисел при доказательстве того, что на отрезке $[0, 1]$ существуют вещественные числа, которые не принадлежать подмножеству $C[0, 1]$ или $U[0, 1]$. Для наглядности, сделаем небольшое изменение множеств $C[0, 1]$ и $U[0, 1]$, а именно поменяем два первых элемента местами, получим новые множества количество элементов, которых не отличается от количества элементов в исходных множествах:
\begin{align}
\begin{split}
&0.100000000\dots\\
&0.000000000\dots\\
&0.010000000\dots\\
&0.110000000\dots\\
&\dots
\end{split}
\end{align}

Применим метод Кантора:
\begin{align}
\begin{split}
&0.\textbf{0}0000000\dots\\
&0.0\textbf{1}000000\dots\\
&0.01\textbf{1}00000\dots\\
&0.110\textbf{1}0000\dots\\
&\dots
\end{split}
\end{align}

Легко видеть, что получается число $0.01(1)$. Однако мы договорились не использовать запись с бесконечным количеством $1$ в двоичной записи, а использовать тождественную запись с цифрами $0$, то есть $0.1$, но такое число уже есть в построенных множествах. Пришли к противоречию. Если бы какое-то из множеств $C[0, 1]$ или $U[0, 1]$ было бы способом нумерации действительных чисел, то число $0.1$ было бы уже посчитано. Если мы восстановим исходный порядок элементов, то картина немного усложнится, диагональный элемент будет $0.1(1) = 1$, $1 \in U[0, 1]$, здесь противоречие сохраняется, но $1 \notin C[0, 1]$. Однако, добавление одного элемента к $C[0, 1]$ не способно изменить его мощность, а значит мы можем прийти к противоречию рассмотрев множество $1 \bigcup  C[0, 1]$, построение диагонального элемента для которого не будет отличаться, поскольку $1 = 1.0(0)$ содержит только $0$ в бесконечной двоичной записи. Очевидно, что существует сколь угодно много последовательностей действительных чисел имеющих условие, что все единицы в двоичной записи числа расположены на позиции меньше диагональной, то есть на диагональной позиции и последующих позициях всегда расположены $0$:
\begin{align}
\begin{split}
&0.\textbf{000000000}\dots\\
&0.x\textbf{00000000}\dots\\
&0.xx\textbf{0000000}\dots\\
&0.xxx\textbf{000000}\dots\\
&\dots
\end{split}
\end{align}

Но для таких последовательностей диагональный метод Кантора либо сразу приводит к противоречию, либо после незначительного изменения порядка элементов, также может потребоваться добавление конечного количества элементов. Данный пример иллюстрирующий неуниверсальность диагонального метода хорошо известен и современные доказательства его легко обходят, запрещая использовать цифры $0$ и $9$ в качестве замены в десятичной записи чисел \cite{matan-diagonal}. Действительно, в двоичной системе можно использовать аналогичный приём, чтобы избежать появления чисел $0.1(1)$ или $0.01(1)$. При построении диагонального элемента мы можем взять не одну цифру, а набор идущих подряд по диагонали цифр произвольной длины, достаточно чтобы комбинация заменяющая эти цифры отличалась от всех комбинаций в строках квадрата, диагональю которого является выбранный набор цифр. Например, будем использовать рамку $2 * 2$ и построим число $b \in [0, 1]$ следующим образом: $b = 0,b_1^1b_2^1b_1^2b_2^2b_1^3 b_2^3\dots$, где $b_1^ib_2^i \neq a^{2i - 1}_{2i - 1}a^{2i - 1}_{2i}$ и $b_1^ib_2^i \neq a^{2i}_{2i - 1}a^{2i}_{2i}$  – комбинация из двух цифр, отличная от комбинаций в двух соответствующих рядах. Всего комбинаций из двух цифр $4$, а значит мы всегда можем выбрать комбинацию отличную от $2$ заданных. Это тот же диагональный метод, только строящийся элемент отличается от каждого элемента множества не в одной конкретной цифре, а в одной из двух. В таком случае для нашего множества по диагонали мы можем получить число $0.01(01)$, $0.01(01) \in U[0, 1]$ (противоречие сохраняется), но $0.01(01) \notin C[0, 1]$. Более того, всякий раз мы можем выбирать произвольную конечную рамку $l * l$, внутри которой выбирать комбинацию, отличную от всех комбинаций в соответствующем квадрате. Тогда у нас получится последовательность рамок $l_1$, $l_2$, $l_3$, $\dots$, $l_j$, $\dots$, и для всякой строки между $l_1 + l_2 + l_3 + \dots + l_{j - 1}$ и $l_1 + l_2 + l_3 + \dots + l_j$ есть цифра в разряде между $l_1 + l_2 + l_3 + \dots + l_{j - 1}$ и $l_1 + l_2 + l_3 + \dots + l_j$ отличающая новое число от числа соответствующего этой строке. Для всех последовательностей действительных чисел, где справа от диагонали всегда расположены нули, это означает, что в качестве диагонального элемента мы можем построить достаточно произвольное число с бесконечной двоичной записью, главное чтобы не существовало разряда $n \in \mathbb{N}$ после которого, заканчивались бы цифры $1$, ведь цифра $1$ в какой-то позиции $k$, гарантированно отличает новый элемент от $k$ первых элементов множества. Иными словами, применяя диагональный метод мы можем прийти к заключению, что множество к которому он применяется ($U[0, 1]$ или $C[0, 1]$) не содержит элементов с бесконечной двоичной записью, произвольным образом чередующей $0$ и $1$, например, число $\pi/4$.

Однако, такого рода приёмы скрывают суть обнаруженного явления, заключающегося в том, что строя элемент отличный от всех в замкнутой алгебре мы в результате получаем элемент, который уже есть во множестве. Мы останавливаемся на этом так подробно, потому что этот пример иллюстрирует проблему в рассуждениях с кристальной наглядностью. Мы определили классы эквивалентности представляющие вещественные числа, но при этом мы не наложили никаких структурных ограничений на бесконечные ряды их определяющие. Для иллюстрации возьмём конечные ряды и определим классы эквивалентности десятизначных чисел, их всего $2^{10}$ и все они представлены в начальных фрагментах множеств $C[0, 1]$ и $U[0, 1]$. Если теперь так случится, что под мощностью мы будем понимать конечность, то диагональное число всегда будет содержаться в начальных сегментах $C[0, 1]$ и $U[0, 1]$ без каких-либо противоречий. Иными словами, ключевой вопрос заключается в том, позволяет ли взаимно однозначное соответствие различать бесконечное количество строк и столбцов, чтобы избежать попадания их в одну категорию эквивалентности как в случае с конечностью? Но это вопрос замкнутости алгебры в которой определена функция соответствия и наличия в ней неподвижных точек, а не универсальное свойство любых биекций. В бесконечном случае диагональ и даже её огромная окрестность содержат только $0$, что буквально означает, что у нас отсутствует возможность заставить диагональный элемент взаимодействовать со всеми элементами счётного множества, то есть какое-то взаимодействие есть, диагональный элемент не может быть абсолютно произвольным, однако, чем дальше мы продвигаемся, тем меньше это взаимодействие и тем больше свободы в продолжении построения, а когда мы приближаемся к бесконечности взаимодействие стремится к нулю, свобода становится безграничной, а наш выбор произвольным, в результате такого построения мы получаем элемент идентичный одному из уже существующих элементов множества. В более строгом смысле, вещественные числа являются классами эквивалентности, которые могут иметь множественные представления в виде двоичных позиционных последовательностей. Например, записи $0.01(1)$ и $0.1(0)$ представляют одно и то же вещественное число. Менее очевидный пример, следует из слабой интерпретации взаимно однозначного соответствия, количество значащих позиций при построении диагонального элемента растёт линейно, но если мы рассмотрим количество значащих цифр в квадратном сегменте фиксированной длины диагонального элемента, то максимальное количество значащих цифр вошедших в данный квадрат растёт логарифмически. При этом, линейная последовательность всё больше и больше отстаёт от логарифмической, которая совпадает с ней значащими цифрами, но обе последовательности линейная и логарифмическая, в конечном итоге, могут порождать экземпляры одного и того же класса эквивалентности вещественного числа. Здесь речь идёт не столько о скорости сходимости, сколько о количественной характеристике множества позиций чисел, полученных в результате на бесконечном шаге. Классическое определение вещественных чисел через классы эквивалентности, не требует количественной тождественности множеств разрядов, а лишь отсутствие конечной позиции различия, но тогда наличие элемента эквивалентного диагональному в множестве вещественных чисел ничему не противоречит. Если под бесконечной мощностью мы понимаем количество элементов за пределами конечного, то диагональный аргумент теряет свою силу. Осознав это, нам остаётся только дать формальное доказательство того, что в некоторых случаях все диагональные элементы, которые только можно построить, действительно могут уже принадлежать нашему множеству. Рассмотрим сначала более детально что происходит на множестве $C[0, 1]$ и какими свойствами обладает это множество.

Лемма 1: {\sl Для произвольного сколь угодно большого натурального $n$ множество $C[0, 1]$ содержит элементы с количеством значащих разрядов равным $n$}

Доказательство:

Предположим, что это не так и существует число $n_{max} \in \mathbb{N}$ такое, что в $C[0, 1]$ все элементы имеют $0$ в любой позиции превосходящей $n_{max}$. Рассмотрим множество $C_{n_{max}}[0, 1]$, в нем содержится элемент $0$, добавленный на нулевом шаге, поскольку $n_{max} + 1 \in \mathbb{N}$, то по построению $C_{n_{max}+1}[0, 1]$ содержит элемент $0.(0\dots0)_{n_{max}}1$, полученный из $0$ и содержащий $1$ в разряде $n_{max} + 1$. Но поскольку $C[0, 1] = \bigcup\limits_{i \in \mathbb{N}} C_i[0, 1]$, то $C[0, 1]$ также содержит этот элемент, что противоречит предположению. Лемма доказана.

То есть, в соответствии с нашим соглашением о потенциальной бесконечности, длина последовательностей значащих разрядов вещественных чисел во множестве $C[0, 1]$ является потенциально бесконечной. Это чрезвычайно важный момент, судя по всему ответственный за то, что мы так долго закрывали глаза на неоднозначности в интерпретации рассуждений Кантора. Когда мы говорим о бесконечности натурального ряда мы одновременно говорим и о бесконечности позиционной записи. Повсеместное использование конечных записей в практике скрывает потенциально бесконечный процесс их построения, что может вводить в заблуждение при анализе бесконечных множеств. Позиционная запись конечна, но она конечна в том же смысле, в каком конечно любое натуральное число и ровно в том же смысле она бесконечна как натуральные числа. С современных позиций этот факт отчётливо виден, позиционная запись определена на уровне мета-теории, в то время как числа являются объектами самой теории, когда появляется бесконечное число, возникает конфликт между теорией и мета-теорией, которая не может отрицать его существование, но и выразить его конечной позиционной записью она не может. Но тогда количество разрядов позиционной записи такого числа не может быть конечным. 

Теорема 1: {\sl для любого $n \in \mathbb{N}$ $C[0, 1]$ содержит все числа $[0, 1)$, количество значащих разрядов (содержащих $1$) которых не превосходит $n$}

Доказательство:

В силу Леммы 1 количество значащих разрядов в вещественных числах принадлежащих множеству $C[0, 1]$ ничем не ограничено. Докажем по индукции, что для любого $n \in \mathbb{N}$, $C_n[0, 1]$ содержит все вещественные числа имеющие $0$ во всех разрядах превосходящих $n$. $C_0[0, 1]$ содержит единственный элемент $0.0(0)$. Если $C_i[0, 1]$ содержит все вещественные числа, количество значащих цифр которых не превосходит $i$. То $C_{i + 1}[0, 1]$ содержит все вещественные числа, количество значащих цифр которых не превосходит $i + 1$. Предположим, что это не так и существует $r$, такое что $r_{i + 1 + k} = 0,  \forall k > 0$ и $r \notin C_{i + 1}[0, 1]$. По построению $C_{i + 1}[0, 1]$ состоит из всех элементов $C_i[0, 1]$, к которым добавлено слагаемое $2^{- i - 1}$, и всех элементов $C_i[0, 1]$ без изменений. Возможно два варианта: $r_{i + 1} = 0$ и $r_{i + 1} = 1$. Если $r_{i + 1} = 0$, то $r \in C_i[0, 1]$, так как $r_{i + k} = 0,  \forall k > 0$, а значит по построению $r \in C_{i + 1}[0, 1]$ и данный случай противоречит предположению, остаётся только случай $r_{i + 1} = 1$. Каким бы ни было число $r$ если в двоичной записи оно содержит $1$ в позиции $i + 1$ за которой следуют все $0$, то число $r - 2^{- i - 1}$ неизбежно содержит все $0$ начиная с позиции $i + 1$. Но тогда, $r - 2^{- i - 1}$ не может принадлежать $C_i[0, 1]$, ведь иначе на итерации $i + 1$ мы получили бы $r - 2^{- i - 1} + 2^{- i - 1} = r \in C_{i + 1}[0, 1]$. Но это противоречит условию, что $C_i[0, 1]$ содержит все вещественные числа, количество значащих цифр которых $i$. Значит $C_n[0, 1]$ содержит все вещественные числа количество значащих разрядов которых не превосходит $n$. Но поскольку $C[0, 1] = \bigcup\limits_{i \in \mathbb{N}} C_i[0, 1]$, то это утверждение справедливо и для $C[0, 1]$. Теорема доказана.

Рассмотрим теперь какие диагональные элементы мы можем построить на множестве $C[0, 1]$. Поскольку это множество является зеркальным отражением множества натуральных чисел в двоичной записи (множества биективно отображаются друг на друга), то каждому элементу множества взаимно однозначно соответствует натуральное число, зеркальным отражением которого этот элемент является. Но тогда по построению Кантора, применённому к потенциально бесконечному множеству $C[0, 1]$, длина $l$ диагонального элемента не может не быть натуральным числом, ведь если $l \notin \mathbb{N}$, то не существует элемента $r$ в $C[0, 1]$ различие с которым определяется разрядом $l$, так как данный элемент должен был бы быть зеркальным (биективным) отображением $l$, которого нет в множестве $\mathbb{N}$, что противоречит определению диагонального построения. Но длина не может быть и натуральным числом, поскольку для любого натурального числа, есть конечное подмножество $C[0, 1]$, которое в силу теоремы 1 содержит этот диагональный элемент. Таким образом, мы можем утверждать, что в потенциально бесконечном случае диагональное построение невозможно, оно бы требовало некоторого нового объекта для описания количества позиций за пределами исходного домена. Также корректным будет утверждение, интерпретирующее потенциальную бесконечность как процесс, что построение такого диагонального элемента никогда не завершиться. Действительно, поскольку множество $C[0, 1]$ состоит из зеркальных отражений (образов) натуральных чисел, а натуральный ряд никогда не заканчивается, то для любого натурального числа есть следующее натуральное число, а значит для любого знака выбранного для диагонального элемента есть следующий знак. Это не является проблемой в теории множеств, поскольку мы явно постулируем существование объектов вне исходного домена, но в потенциально бесконечном случае построение неконсистентно, поскольку у нас отсутствует понятие всех натуральных чисел и выражающий его объект.

Стоит немного углубиться в данную проблему, теорема о несчётности множества вещественных чисел утверждает, что диагональный элемент построен и существует, в то время как построение множества $C[0, 1]$ может продолжаться сколь угодно долго и никогда не будет закончено, ввиду того, что потенциально бесконечная процедура не предусматривает никакого завершения. Чтобы разобраться как такое может быть, давайте рассмотрим последовательность диагональных элементов $\{d_n\}$ построенную на множестве $C[0, 1]$ следующим образом. Начинаем с первого элемента множества и выбираем цифру отличную от первой цифры $c_1$, получаем число $d_1 = 0.d^1$, имеющее ровно один значащий разряд. Следующий диагональный элемент строим из предыдущего, продолжая его двоичную запись цифрой отличной от соответствующей цифры элемента $c_2$: $d_2 = 0.d^1d^2$, количество значащих разрядов - $2$. Будем продолжать данную процедуру неограниченно долго, имея $d_{i - 1}$ построенные по первым $i - 1$ элементам множества $C[0, 1]$, из следующего элемента $c_i$ получаем новый диагональный элемент $d_i$, отличающийся от $c_i$ цифрой в $i$-той позиции, количество значащих разрядов равно его порядковому номеру $i$. Заметим, что такое построение устанавливает взаимно однозначное соответствие между элементами множества $C[0, 1]$, диагональными элементами и натуральными числами, соответствующими номерам элементов и количеству значащих разрядов диагонального элемента, то есть $c_i \Leftrightarrow d_i \Leftrightarrow i$, $\forall i \in \mathbb{N}$. Схожим образом мы можем построить последовательность $\{d_n\}$ для случая диагонального метода использующего рамки произвольной длины, достаточно взять секцию диагонального элемента соответствующую текущей рамке и использовать каждую её цифру для продолжения двоичной записи предыдущего диагонального элемента, включая при этом каждый новый элемент в последовательность диагональных элементов $\{d_n\}$.

Вещественное число, которое должно получиться в результате применения диагонального метода, обозначим как $d$. Все новые элементы последовательности $\{d_n\}$ получаются из старых добавлением новых значащих разрядов, то есть каждый следующий элемент превосходит предыдущий, что значит последовательность $\{d_n\}$ монотонно возрастает, а так как она ограничена отрезком $[0, 1]$, то в соответствии с теоремой Вейерштрасса о пределе монотонной последовательности \cite{matan-weierstrass} она имеет предел, обозначим этот предел как $d_{lim} = \lim_{n\to\infty} d_n$. Легко видеть, что $d$ не может быть ни больше, ни меньше $d_{lim}$, ведь в противном случае согласно правилам сравнения вещественных чисел \cite{matan-decimal} должен существовать конечный разряд $k$, в котором они различны, что исключается процедурами построения обоих чисел, следовательно $d = d_{lim}$. Очевидно, что не выходя за рамки потенциальной бесконечности построить предельный элемент невозможно. В самом деле, предположим противное, что данный предел принадлежит последовательности $\{d_n\}$, то есть существует некоторый порядковый номер $k$, такой что $d_{lim} = d_k$. Но тогда, в силу того, что последовательность монотонно возрастающая, все последующие элементы $d_{k + m}, \forall m > k$ также должны быть равны $d_{lim}$, ведь если есть какой-то элемент, например $d_{k + l}$, больший $d_{lim}$, то бесконечное количество элементов последовательности лежит за пределами окрестности $(d_{k + l} - d_{lim})/2$, что противоречит определению предела \cite{matan-define-limit}. Итак, мы формально доказали, что невозможно построить диагональный элемент, гарантированно отличный от всех элементов потенциально бесконечного множества. В потенциальной бесконечности $d$ - это предел последовательности ${d_n}$, а не её элемент, что исключает его явное построение.

Хорошо, на множестве $C[0, 1]$ явным образом требуемый диагональный элемент не построить, но может быть мы каким-то способом можем его вычислить из последовательности $\{d_n\}$, то есть однозначно сопоставить этой последовательности, какой бы она ни была, некоторое вещественное число $d_{lim}$. Однако, как не сложно убедиться, этот подход также не оставляет никакой возможности сделать заключение о неравномощности множества вещественных чисел на отрезке $[0, 1]$ и множества натуральных чисел. Действительно, множество всех монотонных последовательностей $\{d_n\}$, которые можно построить диагональным методом на счётном множестве, оказывается счётным, так как для любого элемента $d_k$ есть всего два возможных случая: добавление цифры $0$ или $1$ в $k$-ом разряде, то есть с каждым новым разрядом мы просто умножаем количество возможных вариантов диагональных последовательностей на два, что вместе с требованием к каждому элементу иметь конечный номер (сколь угодно большой, но конечный), который, как мы выяснили, взаимно однозначно соответствует количеству значащих разрядов, позволяет нам легко пронумеровать все такие последовательности:
\begin{align}
\begin{split}
0 &\Leftrightarrow \{0.000, 0.000, 0.000, \dots, 0.000, \dots\}\\
1 &\Leftrightarrow \{0.100, 0.100, 0.100, \dots, 0.100, \dots\}\\
2 &\Leftrightarrow \{0.000, 0.010, 0.010, \dots, 0.010, \dots\}\\
3 &\Leftrightarrow \{0.100, 0.110, 0.110, \dots, 0.110, \dots\}\\
\dots &\Leftrightarrow \dots
\end{split}
\end{align}

То есть процедура построения диагональных последовательностей, существенно ограничена и может производить только монотонные последовательности, а количество таких последовательностей для исходного счётного множества, оказывается счётным. Всякая последовательность ${d_n}$ определяется конечным выбором $0$ или $1$ на каждом шаге, что даёт счётное множество (которое условно можно обозначить как $2^{\mathbb{N}}$). Заметим, что считать длину диагональных последовательностей конечно возрастающей не совсем корректно в силу нашего соглашения о потенциальной бесконечности, так как мы не можем указать максимальный номер элемента после которого все последовательности перестают возрастать.

Если счётно множество возможных диагональных последовательностей, то счётно и множество пределов, которые мы можем им сопоставить, по крайней мере, до тех пор пока мы требуем чтобы одной последовательности соответствовал один предел, обозначим множество таких пределов как $L[0, 1]$. Случай множественности пределов соответствует нашей исходной вере, в то что мощность вещественных чисел больше мощности натуральных, однако нам требуется это доказать и построить соответствующий элемент. Таким образом, в силу счётности множества $C[0, 1]$, счётным будет и объединение множеств $C[0, 1] \bigcup L[0, 1]$ \cite{matan-countable-union}, откуда однозначно следует, что, в рамках концепции потенциальной бесконечности, с помощью диагонального метода мы не в состоянии указать ни одного вещественного числа, которое бы не принадлежало счётному множеству $C[0, 1] \bigcup L[0, 1]$. В самом деле, в силу свойства самоподобия бесконечных множеств, которое можно выразить как то, что бесконечное подмножество счётного множества счётно и между ними есть взаимно однозначное соответствие \cite{matan-countable-subset} (отметим, что выбранная формулировка свойства самоподобия избегает выбора между сильной и слабой интерпретациями взаимно однозначного соответствия), объединение множества $C[0, 1]$ с множеством всех пределов $L[0, 1]$, соответствующих диагональным последовательностям, которые могут быть построены на счётном множестве, никак не повлияло на процедуру построения этих диагональных последовательностей. Иными словами, любая последовательность $\{d_n\}$ построенная на множестве $C[0, 1] \bigcup L[0, 1]$ будет принадлежать множеству диагональных последовательностей построенных на множестве $C[0, 1]$ (и на любом другом счётном множестве всех двоичных комбинаций), а значит и её предел окажется во множестве $L[0, 1]$. Иными словами, из-за свойства самоподобия бесконечных множеств диагональный элемент полностью утрачивает связь с множеством, на котором он строится, что позволяет нам объединять его с какими угодно счётными множествами, в частности с теми, которые содержат пределы диагональных последовательностей.

Итак, мы со всей строгостью доказали невозможность диагонального построения на счётном потенциально бесконечном множестве, то есть оперируя только потенциальной бесконечностью и понятием предела мы никоим образом не можем сделать утверждение о существовании диагонального элемента отличного от всех элементов $C[0, 1]$ или $C[0, 1] \bigcup L[0, 1]$. Единственное корректное утверждение, которое мы можем сформулировать имея лишь потенциальную бесконечность выглядит примерно так: для любого сколь угодно большого $n \in \mathbb{N}$ во множестве $C[0, 1]$ может быть построен диагональный элемент с количеством значащих разрядов не превосходящим $n$ отличный от $n$ первых элементов $C[0, 1]$. Однако, $2^n \in \mathbb{N}$ первых элементов $C[0, 1]$ содержат этот диагональный элемент. Или, любому $n$ соответствует диагональный элемент отличный от $n$ первых элементов, но каждой такой паре соответствуют $2^n$ элементов, среди которых есть и тот самый диагональный элемент. Допустимо также говорить о том, что множество натуральных чисел $\mathbb{N}$ равномощно множеству длин диагональных элементов, которые можно построить на множестве $C[0, 1]$, но в то же самое время оба этих множества равномощны множеству степеней двойки $2^n$ и множеству $\{C_n[0, 1]\}$. Более того, этот результат справедлив как для сильной интерпретации взаимно однозначного соответствия так и для слабой, это связано с тем, что домен натуральных чисел замкнут и любая функция здесь представляет собой также потенциально бесконечных объект, комбинаторно представимый внутри домена. Подчеркнём важность данного наблюдения и обратив внимание на то, что в современных теориях множеств мы связываем потенциально бесконечную замкнутую алгебраическую структуру с множеством $\omega$, а не с классом эквивалентности соответствующим кардинальному числу $\aleph_0$, как это сделал Кантор. 

Ни одна из указанных корректных формулировок не позволяет нам сделать вывод о том, что мощность множества вещественных чисел превосходит мощность чисел натуральных. Это не означает, что множество вещественных чисел счётно, ведь числа с бесконечной двоичной записью точно могут не принадлежать $C[0, 1]$ и у нас нет никаких гарантий того, что все такие числа попали в $L[0, 1]$. Однако это позволяет нам однозначно утверждать, что диагональный метод Кантора некорректно использовать для доказательства несчётности множества вещественных чисел, предполагая взаимно однозначное соответствие с некоторым потенциально бесконечным множеством. Авторы полагают, что случай потенциальной бесконечности более менее очевиден и не требовал такой детализации, в конце концов, теория множеств постулирует существование актуальной бесконечности, однако, многочисленные дискуссии в математическом сообществе сформировали у авторов впечатление, что существует определённое недопонимание границ корректного применения диагонального метода, в связи с чем необходимо разобрать все возможности. В действительности, как будет показано далее, определяющим фактором выступает не потенциальность или актуальность бесконечности, а алгебраическая структура объекта, порождаемая правилами его конструирования, а также её специфические свойства, способствующие редукции описания (интерпретации) на уровне мышления, одни структуры допускают применение диагонального аргумента, в то время как другие - нет. Заканчивая с множеством $C[0, 1]$, докажем ещё одно утверждение, которое мы уже упоминали для некоторых частных случаев без доказательства.

Теорема 2: {\sl Для любого вещественного числа $r \in [0, 1]$ множество $C[0, 1]$ содержит последовательность $\{c^r_n\}$, имеющую своим пределом $r$}

Доказательство:

Каждое вещественное число $r \in [0, 1]$ имеет двоичную запись $r = 0.r_1r_2r_3\dots$, где $r_1$, $r_2$, $r_3$ – цифры 0 или 1. Рассмотрим последовательность $\{c^r_n\}$ полученную следующим способом: $c^r_0 = 0$, $c^r_i = c^r_{i - 1}$, если $r_i = 0$, в противном случае $c^r_i = c^r_{i - 1} + 2^{-i}$, так или иначе $c^r_i \in C_i[0, 1]$. По построению двоичная запись $c^r_i$ совпадает с двоичной записью $r$ до $i$-го разряда:
\begin{equation}
c^r_i = 0.r_1 r_2 \dots r_i
\end{equation}

Но тогда $c^r_i$ есть ничто иное, как нижнее двоичное приближение числа $r$ порядка $i$ и следовательно предел последовательности $\{c^r_n\} = r$ \cite{matan-bottom-limit}. Легко видеть также, что верхняя граница отрезка - число $1$, является пределом последовательности $0.1$, $0.11$, $0.111$, $\dots$, элементы которой принадлежат $C[0, 1]$. Теорема доказана.

Таким образом, с одной стороны пределов не может быть больше чем последовательностей в $C[0, 1]$, то есть чем натуральных чисел, с другой для каждого вещественного числа имеется соответствующая последовательность $C[0, 1]$ и натуральное число. В действительности никакого противоречия в этом нет, поскольку мы используем классическое представление о вещественных числах, где бесконечная последовательность не подразумевает никаких структурных ограничений на количество элементов, в то время как в современная теория различает последовательности различной трансфинитной длины. В контексте диагонального построения можно сформулировать это так: всякое вещественное число, которое можно сопоставить с диагональной последовательностью можно сопоставить и с некоторой другой последовательностью $C[0, 1]$. Итак, мы не смогли добиться какого-либо успеха, используя потенциальную бесконечность, давайте теперь исследуем границы применимости диагонального метода в актуально бесконечном случае. Потенциально бесконечное множество $C[0, 1]$ не содержит пределов последовательностей своих рациональных элементов, из-за чего мы не можем построить такого предела на диагонали. В принципе, нельзя исключать существование последовательностей элементов множества $C[0, 1]$, предел которых невозможно получить диагональным методом, и их количество вполне может оказаться несчётным, но даже если мы придумаем способ получения таких последовательностей, то он совершенно точно не будет диагональным, ведь с помощью последнего можно построить лишь счётное множество пределов. Иными словами, это даст нам некоторое другое доказательство несчётности, которое выходит за рамки обсуждения диагонального метода. Включив в множество $C[0, 1]$ счётное множество каких-то предельных элементов, диагональное построение всё также будет давать счётное множество пределов. Рассмотрим актуально бесконечное множество $U[0, 1]$, возможно оно позволит построить диагональную последовательность предел, которой не будет элементом самого множества $U[0, 1]$, в конце концов, Кантор оперировал именно актуальной бесконечностью координат (хоть и избегал называть такую систему координат последовательностью) \cite{diagonal}:

Теорема 3: {\sl Для любого вещественного числа $r \in [0, 1]$ множество $U[0, 1]$ содержит последовательность $\{u^r_n\}$, имеющую своим пределом r, при этом $u^r_{\infty} = r \in U[0, 1]$}

Доказательство:

Воспользуемся способом построения последовательности $\{u^r_n\}$ из теоремы 2. Согласно определению множеств $U[0, 1]$ и $C[0, 1]$ последовательность $\{u^r_n\}$ содержит все элементы последовательности $\{c^r_n\}$, следовательно предел последовательности $\{u^r_n\} = r$.

Докажем, что $u^r_{\infty}$, то есть элемент множества $U[0, 1]$, такой что для всякого $i \in \mathbb{N}$ $i$-ый разряд $u^r_{\infty}$ совпадает с $i$-ым разрядом $\{u^r_i\}$, является пределом последовательности $\{u^r_i\}$. Согласно определению предела \cite{matan-define-limit} это означает, что для всякого $\epsilon > 0$ существует такой номер $n_{\epsilon}$, что для всех $n > n_{\epsilon}$ выполняется неравенство:
\begin{equation}
|u^r_{\infty} - u^r_n| < \epsilon
\end{equation}

Каким бы ни было число $\epsilon$ согласно принципу Архимеда \cite{matan-archimed} существует натуральное число $k_{\epsilon} > 1/\epsilon$. Поскольку $2^{k_{\epsilon}} = (1 + 1)^{k_{\epsilon}} > 1 + {k_{\epsilon}} + \dots > k_{\epsilon}$, то $2^{k_{\epsilon}} > 1/\epsilon$ и $\epsilon > 2^{-k_{\epsilon}}$. По построению $u^r_{\infty}$ может отличатся от $u^r_i$ только цифрами в разрядах превышающих $i$, значит $u^r_{\infty} - u^r_i \leq 2^{-i}$. Тогда $|u^r_{\infty} - u^r_{k_{\epsilon}}| \leq 2^{-k_{\epsilon}} < \epsilon$, следовательно для всех $k > k_{\epsilon}$
\begin{equation}
|u^r_{\infty} - u^r_k| \leq 2^{-k} < 2^{-k_{\epsilon}} < \epsilon
\end{equation}

То есть последовательность $\{u^r_i\}$ имеет своим пределом $u^r_{\infty}$. Но тогда, $u^r_{\infty} = r$ в силу единственности предела числовой последовательности \cite{matan-single-limit}. Теорема доказана.

Следствием теоремы 3 является, то что все вещественные числа на отрезке $[0, 1]$ принадлежат множеству $U[0, 1]$. А это значит: какое бы число мы ни попытались построить по диагональному методу - оно всё равно уже будет находиться внутри множества $U[0, 1]$. Таким образом, диагональный аргумент сам по себе не даёт никакой структурной информации о соотношении множеств $\mathbb{N}$ и $\mathbb{R}$. Ему необходима поддержка, которая уже фиксирует определённые структурные ограничения в отношении того, что мы называем бесконечностью. Такие ограничения способно дать взаимно однозначное соответствие, однако и оно существенно зависит от той алгебраической структуры, внутри которой определено, а значит, в конечном итоге, от предполагаемой морфологии бесконечных объектов. Действительно, если бесконечность мыслится как класс эквивалентности, то и взаимно однозначное соответствие становится таким классом (слабая интерпретация), поскольку оно является объектом той же алгебры, что и бесконечный объект. Если бесконечность мыслится как фиксированное множество, то и взаимно однозначное соответствие становится фиксированным множественным объектом (сильная интерпретация). Получается всё сводится к тому, какую структуру имеет множество $U[0, 1]$, тому, как именно мы аннотируем системы строк и столбцов и какие операции считаем корректными, что в свою очередь непосредственно связано с предположениями относительно свойств бесконечности. Давайте рассмотрим какие именно предположения мы можем сделать.

\begin{itemize}
    \item Бесконечность универсальна, единственна, бесструктурна, неограниченно рефлексивна. Например, бесконечность - это всё за пределами конечного. Тогда бесконечное количество столбцов неотличимо от бесконечного количества строк, а значит оба множества отображаются друг на друга. Однако, в таком случае в силу теоремы 3 любое вещественное число полученное диагональным методом уже присутствует в множестве $U[0, 1]$. Следовательно диагональное построение либо нерезультативно - не позволяет сделать вывод о противоречии, например, как в случае когда бесконечность является классом эквивалентности с неограниченной рефлексивностью, либо неконсистентно и должно быть исключено из теории, например, при помощи стратификации. В данный класс предположений попадают все конструкции, которые моделируют бесконечность с помощью потенциально бесконечных построений, где как мы установили невозможно сконструировать диагональный элемент. А также все замкнутые алгебры, в частности, ряд современных альтернативных концепций, где предполагается неограниченная рефлексивности бесконечности: $\infty + 1 = \infty + \infty = \infty * \infty = \infty$ (замкнутость относительно всех операций), допускающих счётность вещественных чисел. Здесь диагональное построение можно легко заблокировать, например, указав, что алгоритмически граница недостижимости конечными вычислениями не видит всей матрицы $U[0, 1]$. В частности, движение по диагонали связано с двумя атомарными операциями перехода к следующей строке и следующему столбцу, а значит диагональный элемент уходит за границу конструктивной бесконечности, ещё до того как он может быть построен.   
    \item Бесконечность имеет структуру, ограниченно рефлексивна, множественна. Например, бесконечность - это принадлежность универсальному пространству координат, как предложил Кантор. Такое пространство должно быть рефлексивно, без такого предположения диагонального построения на актуально бесконечном множестве $U[0, 1]$ просто не существует, в силу того, что предел отношения количества столбцов к количеству строк $\lim_{n\to\infty} n/2^n = 0$, то есть между ними не может быть биекции. Тогда чтобы определить диагональный элемент нужно задействовать в его построении все строки, для чего нужно принять сразу два серьёзных предположения.
\begin{enumerate}
    \item Универсальность. Допустимость говорить о взаимно однозначном отображении позиций всех вещественных чисел друг на друга одновременно. Или если подойти с другой стороны, существует некая общая структура (пространство), которая позволяет одновременно и взаимно однозначно индексировать позиции двоичных разрядов всех вещественных чисел из $[0, 1]$ сразу. Нам не достаточно существования некоторого типа (класс эквивалентности) индексирующих множеств, который мы называем счётным. Поскольку из того, что для каждого вещественного числа в данном типе есть индексирующее множество, в общем случае, не следует существование множества пригодного для индексации позиций всех вещественных чисел. Нам нужно чтобы данный тип был, некоторым образом ограничен, например, чтобы существовало объединение множеств всех позиций, которое мы могли бы использовать как эталон, в противном случае, поскольку нет ограничения, то нет и возможности задать диагональный элемент. В таком случае, предполагаемой математической структуры действительно может оказаться недостаточно, чтобы проиндексировать все элементы $U[0, 1]$, то есть речь идёт о предположении, что бесконечность обладает в каком-то смысле общей идентифицируемой инвариантностью.
    \item Тотальность. Ключевое предположение доказательства заключается в существовании биекции (сильная интерпретация) между столбцами и строками, которое и позволяет прийти к противоречию. Нам нужно предположить, что множество индексов столбцов и множество индексов строк - это потенциально сопоставимые друг с другом множества, а не просто множества, связанные отношением эквивалентности выражающим некоторое уникальное свойство бесконечности. Иными словами, нам требуется тотальность структуры, способность теории потенциально описывать каждый её элемент. В частности, тотальная вполне упорядоченность, но не как эмерджентно возникающее свойство выразимых (достижимых) элементов, а как имманентное свойство самой бесконечной целостности. Только такая тотальность позволяет проецировать структуру универсального пространства координат, на композицию векторов данного пространства. Без тотальности диагональный элемент может оказаться просто объектом, который не описывает наша теория, объектом за пределами описательной силы теории.
\end{enumerate}
\end{itemize}

Иными словами, диагональный метод не является инструментом исследования структуры бесконечности, наоборот, предположение о морфологии таких объектов определяет допустимость использования диагонального метода. Само построение диагонального элемента сильно зависит от того, какую именно модель актуальной бесконечности мы используем. В одних интерпретациях для множества $U[0,1]$ диагональная процедура просто не существует. В других же диагональный элемент действительно появляется и позволяет различать количественные характеристики множеств. В любом случае мы приходим к выводу: классический диагональный метод Кантора не является универсально применимым даже на актуально бесконечных множествах. Особенно это заметно, когда доказательство приводится без явного указания предпосылок, обуславливающих его корректность. Кантор в отличии от многих последующих работ явно формулирует ключевые предположения, которые он считает достаточными. Более того, он выделяет две алгебры: кардинальную и порядковую, в которых морфология бесконечных объектов принципиально различна, однако в современной теории мы видим упрощение этой принципиальной позиции. Таким образом, с формальной точки зрения заключение о несчётности множества вещественных чисел не является общезначным и зависит от принятых предположений. Более того, как будет показано в следующи работах, смысл данного заключения различается в рамках теории Кантора и в современной теории. Натуральные и вещественные числа - это первичные объекты, которые мы ощущали и использовали задолго до формализма теории множеств. У нас гораздо больше оснований считать их реально существующими, чем, например, аксиому бесконечности или другие абстрактные конструкции. Именно поэтому неоднозначность, которую мы обнаруживаем при внимательном рассмотрении диагонального аргумента, даёт очень ценную подсказку: именно здесь, в основаниях теории множеств, скорее всего и скрыта проблема, которая остаётся не до конца осмысленной и сегодня.

Перейдём теперь к оригинальной версии доказательства Кантора \cite{cantor-original}. Пусть $m$ и $w$ - два каких-либо исключающих друг друга признака, рассмотрим совокупность $M$ элементов $E = (x_1, x_2, \dots, x_n, \dots)$ зависящих от бесконечно многих координат $x_1, x_2, \dots, x_n, \dots$, где каждая из этих координат есть $m$ или $w$. Пусть $M$ - совокупность всех элементов $E$. Элементами совокупности $M$ являются, например, следующие:
\begin{align}
\begin{split}
E^1 &= (m, m, m, m, \dots)\\
E^2 &= (w, w, w, w, \dots)\\
E^3 &= (m, w, m, w, \dots)\\
\end{split}
\end{align}

Теорема: {\sl Если $E_1, E_2, \dots, E_n, \dots$ - какая-либо просто бесконечная последовательность элементов многообразия $M$, то всегда существует такой элемент $E_0$ многообразия $M$, который не совпадает ни с каким $E_n$.}

Доказательство:

Пусть
\begin{align}
\begin{split}
\\[-2\baselineskip]
E_1 &= (a_{11}, a_{12}, \dots, a_{1n}, \dots)\\
E_2 &= (a_{21}, a_{22}, \dots, a_{2n}, \dots)\\
\dots &= \dots\\
E_m &= (a_{m1}, a_{m2}, \dots, a_{mn}, \dots)\\
\dots &= \dots
\end{split}
\end{align}

Здесь $a_{mn}$ являются признаками $m$ или $w$. Определим теперь последовательность $b_1, b_2, \dots, b_n, \dots$ так, чтобы $b_n$ был тоже равен только $m$ или $w$ и отличен от $a_{nn}$. Итак, если $a_{nn} = m$, то $b_n = w$, а если $a_{nn} = w$, то $b_n = m$.

Если теперь мы рассмотрим элемент $E_0 = (b_1, b_2, \dots, b_n, \dots)$ многообразия $M$, то очевидно, что равенство $E_0 = E_k$ не может иметь места ни для какого положительного целочисленного значения $k$, так как в противном случае для соответствующего $k$ и для всех целочисленных значений $n$ было бы $b_n = a_{kn}$, а значит, в частности, $b_k = a_{kk}$ что исключается определением $b_k$. Теорема доказана.

Из этой теоремы непосредственно следует, что совокупность всех элементов многообразия $M$ нельзя представить в форме последовательности $E_1, E_2, \dots, E_n, \dots$, так как в противном случае мы столкнулись бы с противоречием, что некая вещь $E_0$ как была бы, так и не была бы элементом многообразия $M$.

Кантор формулирует теорему максимально абстрактно, на множестве элементов зависящих от бесконечного набора координат, с неопределённым набором свойств. Примечательно, что Кантор вводит систему координат, усиливая требования к наборам признаков, ведь ему требуется возможность одновременной индексации координат во всех элементах совокупности $E_1, E_2, \dots, E_n, \dots$ - свойство, которое было нетривиально доказать для последовательностей, по крайней мере, во времена Кантора. С современных позиций можно говорить о наборе координат как о фиксированном счётном множестве или универсальном (общем) двухзначном пространстве, в то время как под последовательностью, в данном контексте, нужно понимать некоторый объект принадлежащий классу эквивалентности тотально упорядоченных совокупностей. Иными словами, Кантор оставляет возможность того, что натуральные числа - это не множество, а класс эквивалентности подобных друг другу множеств. По всей видимости, это отнюдь не случайно, и связано с поиском корректного языка для описания открытых им объектов, который окончательно формируется в его следующих работах. Кантор, судя по всему, вполне осознанно, фиксирует ключевое свойство координатной системы - универсальность (все объекты описываются единой системой координат), и одновременно ключевое свойство последовательности - тотальную упорядоченность. То есть вместо потенциально бесконечных последовательностей, мы рассматриваем полностью упорядоченную законченную бесконечную структуру. Стоит обратить внимание на философскую глубину доказательства, изолирующего исключительно необходимые свойства объектов, что существенно облегчает анализ области применимости, чем часто пренебрегают. Если переводить доказательство на современный формальный язык, то вывод оказывается потрясающе чистым: из предположения существования универсальной системы координат представимой в форме тотально упорядоченной бесконечной последовательности (эталона), следует невозможность представления таким эталоном совокупности всех точек данного двухзначного пространства. Из-за отсутствия во времена Кантора формальных требований, многие предположения нуждающиеся в фиксации вводятся через непосредственное использование. В частности, из доказательства можно сделать вывод, что существенным свойством элементов многообразия $M$ является процедура сравнения их между собой, следующая из упорядоченной структуры координат. Кантор, определяет равенство элементов как равенство всех координат одного элемента всем координатам другого в каждой позиции $n$ от 1 до бесконечности. Иными словами, $E_n = E_m$, тогда и только тогда, когда $a_{nk} = a_{mk}$ для всех $k = 1, 2, \dots$ (нам важно, чтобы счёт мог проиндексировать всё множество бесконечных координат). То есть в момент сравнения мы последовательно сопоставляем координаты одного элемента другому и если находим конечную позицию, где координаты не совпадают, то элементы не равны, если же такая позиция нам никогда не встретится, то элементы равны. Обратим внимание ещё на один любопытный момент, мы используем невычислимость различия элементов $E$ и вещественных чисел как факт равенства или даже его определение \cite{compare-undecidable}. 

Кантор строит некоторую последовательность координат и предлагает рассмотреть такой элемент, координаты которого совпадают с этой последовательностью. То есть, смысл доказательства состоит в том, что взяв любое фиксированное бесконечное множество координат, мы можем определить множество $M$ или пространство, для произвольной последовательности элементов которого, можно указать некоторый элемент $E_{0}$, координаты которого отличаются от координат всех элементов этой последовательности. Полностью упорядоченная универсальная бесконечная система координатных позиций позволяет определить бесконечное многообразие $M$, которое будучи полностью упорядочено не принадлежит тому же классу эквивалентности, требуя для упорядочения индексов, выходящих за пределы системы координатных позиций. Иными словами, два выделенных нами свойства бесконечных объектов: универсальность системы координат и тотальный порядок её элементов, приводят к разделению классов эквивалентности (типов) в отношении возможности взаимно однозначного соответствия в рамках алгебраической структуры пространства. Если мы связываем определение последовательности с такой структурой, то координаты допускают представление в форме просто бесконечных последовательностей, в то время как многообразие $M$ не допускает. Здесь Кантор уже невероятно близок к разделению кардинальной и порядковой арифметики, он высказывает идею, что мощности представляют собой кардинальные числа, в то время, как в предыдущих своих работах числами Кантор называл только порядковые типы, тем не менее, независимые алгебры ещё не выкристаллизовались и мощность ещё понимается как общезначное свойство невозможности взаимно однозначного соответствия, в чём он усомнится впоследствии и введёт два вида такого соответствия. Однако, аксиоматические теории множеств, судя по всему, наследуют незавершенную упрощённую интерпретацию, в рамках которой предполагается универсальность понятия числа, означающая количественную тождественность множества колонок и строк, то есть множество номеров колонок тождественно множеству номеров строк. Хотя сам Кантор, подобных утверждений нигде не делает, и даже более того, как авторы намереваются показать, он не считает множества одной мощности количественно равными, а при равенстве конечных элементов тождественными, поэтому приписывать ему современную интерпретацию было бы совершенно неправильно. 

Точность формулировок свидетельствует о том, что Кантор сознательно использует именно актуально бесконечные объекты имеющие упорядоченную структуру, а не потенциально бесконечные объекты - последовательности, в классическом смысле. В случае последних, множество таких элементов, как $E_0$, может быть счётным, этот факт связан с важным для Кантора свойством бесконечных множеств быть подобными своей части, что мы наблюдали на примере множества $C[0, 1]$, в частности, этот факт уже присутствует в доказательстве счётности множества рациональных чисел. В случае потенциальной бесконечности нет универсального ограничения на класс эквивалентности, поэтому когда мы строим элемент отличный от всех элементов бесконечного множества, мы можем в результате получить элемент, для которого во множестве уже присутствуют эквивалентные элементы. Это указывает нам на область применимости доказательства, непосредственно связанную с тем, как мы определяем структуру системы бесконечных координат. То есть, случай потенциально бесконечной системы координат, когда все её позиции конечны, полностью совпадает с уже разобранным случаем для вещественных чисел множества $C[0, 1]$, понятие предела нам здесь не требуется, поскольку мы можем просто сопоставить всякой диагональной последовательности некоторый элемент, и этого будет достаточно, поскольку счётность диагональных последовательностей влечёт и счётность сопоставленных элементов. В случае актуально бесконечной система координат, также как и для множества $U[0, 1]$, мы можем строго показать, что бесконечная упорядоченная система $E_1, E_2, \dots, E_n, \dots$ элементов множества $M$ может содержать все элементы этого множества. В самом деле, пусть $M[w, m]$ - такая система элементов многообразия $M$, построенная по тем же принципам, что множества $C[0, 1]$ и $U[0, 1]$, где 0 соответствует $w$, a 1 соответствует $m$.

Теорема 4: {\sl Бесконечная система $M[w, m]$ содержит все элементы многообразия $M$}

Доказательство:

Предположим, что существует некоторое $E_0 = (b_1, b_2, \dots, b_n, \dots)$, такое что равенство $E_0 = E_k$ не имеет места ни для какого положительного целочисленного значения $k$. Используя определение равенства последовательностей неявно присутствующее в доказательстве теоремы Кантора, будем последовательно исключать из рассмотрения все элементы различающиеся с $E_0$ в позициях $n = 1, 2, \dots$. Возможны два случая: сколько бы мы не проводили нашу процедуру, мы никогда не исключим всех элементов, либо существует какой-то последний шаг, на которым мы исключаем последние оставшиеся элементы.

Если исключение всех элементов не произойдёт никогда, значит оставшиеся в рассмотрении элементы равны $E_0$, здесь мы используем свойство невычислимости равенства элементов. В самом деле, возьмём первый из оставшихся элементов $E_x$ если бы он не равнялся $E_0$, то должна была бы существовать некоторая целочисленная положительная позиция $l$, такая что $b_l \neq a_{xl}$. Но в таком случае данный элемент был бы исключён из рассмотрения на шаге $l$. Следовательно, $E_0 = E_x$. Но это противоречит предположению, что существует элемент отсутствующий в $M[w, m]$.

Значит единственный вариант - это то, что в какой-то момент, на шаге $j$, мы исключили из рассмотрения все элементы. Но из этого следует, что все элементы $M[w, m]$ различаются с $E_0$ в позиции не превосходящей $j$. Рассмотрим $M_j[w, m]$ соответствующее $C_j[0, 1]$ и $U_j[0, 1]$, по построению $M_j[w, m]$ содержит все возможные комбинации признаков $w$ и $m$ длины $j$, но это означает, что $E_0$ совпадает с одним из элементов $M_j[w, m]$, во всех позициях включая $j$, обозначим такой элемент $E^j_0$. Но поскольку $M_j[w, m] \subset M[w, m]$ то $E^j_0 \in M[w, m]$, и мы никак не могли удалить его за $j$ шагов. Следовательно $M[w, m]$ содержит все элементы многообразия $M$. Теорема доказана.

Данное доказательство требует некоторых предположений обосновывающих бесконечную процедуру исключения из рассмотрения, с формальной точки зрения здесь мы опираемся на сформулированную выше аксиому о возможности бесконечных математических процедур, без которой, в частности, невозможно говорить о построении диагонального элемента. В рамках данной работы такое объяснение удовлетворительно, поскольку её цель заключается в том, чтобы продемонстрировать границы применимости диагонального метода. Отметим также, что приведённое выше доказательство не противоречит аксиоматикам теории множеств как таковым, вопрос сводится лишь к определению понятия последовательности и её свойств, а также возможности называть $M[w, m]$ последовательностью. Таким образом, как и в случае с $U[0, 1]$, возможность прийти к противоречию диагонального метода, полностью зависит от исходных предположений относительно природы бесконечного. Отметим также, что заключение Кантора в контексте данной работы совершенно верно, в случае фиксированной универсальной полностью упорядоченной системы координат, биекция между ней и полностью упорядоченной производной системой $M$ невозможна. Невероятно филигранной и строгой в математическом смысле является также обозначенная Кантором связь между данной теоремой и доказанного им ранее факта непредставимости вещественных чисел в форме последовательности. Кантор не утверждает, что вещественные числа можно описывать универсальной тотально упорядоченной системой координат рассматриваемой в доказательстве, он утверждает лишь, что обе теоремы показывают расслоение классов эквивалентности мощности на типы при определённых предположениях о структуре. В теории вещественных чисел мы делаем одни предположения о их структурной организации, в рамках данной теоремы мы делаем, возможно, минимальные, но другие предположения. Сведение вещественных чисел к векторам в универсальном тотально упорядоченном пространстве - это сильное упрощение аксиоматических подходов, возникших впоследствии, а не позиция Кантора. И именно с этим упрощением связано различие континуума в понимании Кантора и континуума в современном контексте, которое проявляется в существенном расхождении того, что Кантор подразумевает под несчётностью континуума и что подразумевают современная теория. Таким образом, формулировки Кантора допускают интерпретацию того, что универсальная тотально упорядоченная система координат является фрагментом теории вещественных чисел, не покрывающим их полностью. Необходимо также отметить, что строгая формулировка свойств универсальности и тотальности обеспечивает корректность теоремы применительно к классам эквивалентности (типам множеств), с которыми работает Кантор, в частности для порядкового типа $\omega$. И вывод Кантора корректно обобщается: при фиксации ограниченного класса эквивалентности упорядоченных множеств, многообразие $M$ не может быть упорядочено в объект принадлежащий тому же классу, даже если мы говорим о слабой интерпретации взаимно однозначного соответствия, поскольку оно как класс эквивалентности множества пар также подпадает под ограничение. Любопытно, что в своей статье Кантор делает ещё одно утверждение, акцентируя внимание на сути возникшего противоречия - неопределённость в принадлежности элемента множеству. В многочисленных воспроизводствах диагонального метода, редко приходится сталкиваться с этим важным моментом, очень точно, тем не менее, характеризующим возникшую ситуацию, которая, в свете обнаруженных впоследствии парадоксов теории множеств, должна была бы вызвать настороженность. 

Для окончательного прояснения, предполагаемой авторами позиции Кантора, приведём здесь одно занимательное наблюдение. Давайте рассмотрим подмножество $M^m$ множества $M$ состоящее из элементов $E^m$, которые либо имеют ровно одну смену признака и начинаются с $m$, либо не имеют ни одной смены признака и начинаются с $m$. В этом множестве, есть элемент $E^m_* = (m, m, ..., m, ...)$, для которого серия признака $m$ никогда не прерывается. Легко видеть, что полученное множество взаимно однозначно соответствует множеству координатных позиций $\mathbb{E}$. Действительно, примитивная форма записи чисел - повторение какого-либо признака (знака, чёрточки) определённое количество раз, соответствующее этому числу. Таким образом, всякой координатной позиции мы можем сопоставить $E^m$, где данная позиция является финальной для признака $m$. Ввиду того, что мы разрешили бесконечные процедуры, бесконечное количество чёрточек должно обозначать некоторое число, которое соответствует элементу $E^m_*$ (современные представления не позволяют описывать такие чисел, так как они противоречат интерпретации некоторых принятых нами предположений, но здесь мы договорились делать минимум предположений о сущности бесконечного). Добавим теперь для симметрии во множество $M^m$ элемент $E_w = (w, w, ..., w, ...)$. Либо, что тоже самое, изменим условие на элементы множества $M$, так чтобы этот элемент попадал в результирующее множество $M^m \bigcup \{E_w\}$, то есть будем выбирать те элементы, которые либо имеют ровно одну смену признака и начинаются с $m$, либо не имеют ни одной смены признака. Очевидно, что есть только один элемент, не имеющий смены признака и не начинающийся на $m$ и это $E_w$. По построению $M^m \subset (M^m \bigcup \{E_w\}) \subset M$. Пусть $E_w, E^m_1, E^m_2, \dots, E^m_n, \dots$ - бесконечная упорядоченная система элементов множества $M^m \bigcup \{E_w\}$, такая что все элементы в ней расположены в порядке возрастания длины серии признака $m$. То есть:
\begin{align}
\begin{split}
E_w &= (w, w, w, \dots, w, \dots)\\
E^m_1 &= (m, w, w, \dots, w, \dots)\\
E^m_2 &= (m, m, w, \dots, w, \dots)\\
\dots &= \dots\\
E^m_n &= (m, m, m, \dots, a_{n+1} = w, \dots)\\
\dots &= \dots
\end{split}
\end{align}

Если вдруг по какой-то причине мы не можем предъявлять требования к взаимно однозначному соответствию относительно определённого порядка элементов, хотя сам Кантор считал их абсолютно законными \cite{cantor-relation}, то ровно по этой же причине мы не можем и исключить возможность определённого порядка элементов, поскольку требование отсутствия определённого порядка, есть ничто иное как требование некоторого другого порядка. Однако, в таком случае, используя процедуру предложенную Кантором мы получаем элемент $E_0 = E^m_* = (m, m, \dots, m, \dots)$ (предположение, что какая-то координата не имеет признак $m$ противоречиво), который согласно построению принадлежит множеству $M^m \bigcup \{E_w\}$, но в силу того, что он является диагональным элементом, он не может принадлежать этому множеству. Следовательно, согласно предположению о диагональном методе, упорядоченная система $E_w, E^m_1, E^m_2, \dots, E^m_n, \dots$ не может взаимно однозначно отображаться на упорядоченную систему $E^m_1, E^m_2, \dots, E^m_n, \dots$ с сохранением порядка. Таким образом, диагональный аргумент иллюстрирует ограничение порядковой арифметики Кантора, требующее аккуратного определения структуры порядковых типов. С позиций теории множеств мы бы сказали, что $E^m_1, E^m_2, \dots, E^m_n, \dots, E^m_*$ уже имеет тип $\omega + 1$, а диагональный элемент $E_0 = (w, w, \dots, w, \dots)$ указывает, что в части без $E^m_*$ уже задействованы все позиции. Данное наблюдение обнажает простой факт того, что структура вполне упорядоченности не является единственной упорядоченной морфологией бесконечности, именно поэтому Кантор в своих исследованиях говорит о типах множеств (классах эквивалентности и классах подобия), а не о самих множествах. Однако, структура типов позволяет математически строго выделить свойство подобия собственной части, отличающее конечное от бесконечного, которое имело ключевое значение для Кантора. Он не принимал альтернативную морфологию, наделяющую бесконечные числа набором свойств, идентичных конечным, поскольку видел порочный круг в императивном разграничении объектов с одинаковым набором свойств на конечные и бесконечные. Наивные попытки утверждать, будто однотипные объекты можно классифицировать на конечные и бесконечные на основании одной интуиции - когда бесконечное возникает как ''не конечное'', а конечное имеет лишь смысл ''не бесконечного'', появляясь в едином акте дифференциации, обосновать который можно только через него самого, - Кантор справедливо не считал математичными. Иными словами, вопрос бесконечной морфологии для Кантора сводился к достаточному математическому обоснованию, которого он не смог найти для структуры самих множеств, но нашёл для типов.

Ещё раз подчеркнём, что все рассуждения о многообразии $M$, изложенные выше, оперируют исключительно актуально бесконечными множествами, в рамках простой аксиомы декларирующей их формальную законность. Именно в таких множествах возникают парадоксы связанные с тем, что элемент одновременно принадлежит множеству и не принадлежит ему, это исключительно важный вопрос связанный с перспективой нашего описания, в одних описаниях верно первое, в других - второе. В классическом представлении потенциально бесконечных величин, безраздельно властвует принцип математической индукции, который не допускает парадоксов, что мы наблюдали при попытке применить диагональный метод ко множеству $C[0, 1]$. Точно также, не возникает парадокса со множеством $M^m \bigcup \{E_w\}$, где легко доказать по индукции, что диагональный элемент множества $M^m \bigcup \{E_w\}$ равен последнему элементу этого множества, для любой сколь угодно большой длины $n$. Действительно, на множестве последовательностей длины $1$, множество $M^m \bigcup \{E_w\}$ состоит из двух элементов $\{w\}$ и $\{m\}$, диагональный элемент равен $\{m\}$ тождественнен последнему элементу, на множестве последовательностей длины 2, в $M^m \bigcup \{E_w\}$ три элемента: $\{w, w\}$, $\{m, w\}$ и $\{m, m\}$, диагональный элемент - $\{m, m\}$ и т. д.

Таким образом, $M[w, m]$ может быть и не является бесконечной последовательностью, однако, доказать это при помощи диагонального метода, не сделав дополнительных предположений относительно того, что именно мы называем бесконечной последовательностью (например, подразумеваем ли мы также трансфинитные последовательности), мы не в состоянии. В то же время, стоит заметить, что переформулировка задачи следующим образом: существует ли диагональный элемент отличный от любого элемента множества $M[w, m]$ или корректна (существует, определена) ли процедура построения такого элемента, - совместима с предположением, что $M[w, m]$ - бесконечная последовательность. Такая задача в зависимости от наших предположений может иметь решение, устанавливающее, что такого диагонального элемента не существует, равно как и процедуры его построения. В частности, для множества $M[w, m]$ мы можем заметим, что всегда вынуждены выбирать признак $m$, поскольку все элементы на диагонали имеют признак $w$. Но тогда, мы получим элемент $E^m_* = (m, m, ..., m, ...)$, который уже есть во множестве, что противоречит условию задачи. Какие бы ухищрения мы не использовали при построении диагонального элемента, множество $M[w, m]$ не оставляет ни одного кандидата для диагонального элемента. Еще раз отметим также, что Кантор в своей статье нащупывает корректный язык для описания обнаруженных им явлений, он уже понимает, что его интуитивное универсальное понятие мощности не совсем подходит для описания различных проявлений взаимосвязей между множествами, язык, который он найдёт и изложит в своей кульминационной работе. Авторы позволят себе упомянуть, что приведённые выше рассуждения, несмотря на то, что они могут вызывать некоторый диссонанс у специалистов в теории множеств, не просто не противоречат аксиоматическим теориям множеств, а являются математически корректными конструкциями имеющими достаточное обоснование, но к сожалению невозможно объять необъятное и сократить предстоящий сложный путь. 

Давайте посмотрим на проблему с другой стороны - через понятие взаимно однозначного соответствия. Согласно Кантору, такое соответствие - это всего лишь способ разделить множества на некоторые структурные классы эквивалентных множеств. Однако, в рамках современных теорий мы используем сильную интерпретацию, превращающую структурную эквивалентность в универсальную количественную тождественность. По аксиоме экстенсиональности два множества равны, если у них один и тот же набор элементов. В этом случае между ними существует биекция, где каждый элемент соответствует самому себе. Но обратное, вообще говоря, не верно. Например, множество чётных натуральных чисел и множество удвоенных натуральных чисел находятся в очевидном взаимно однозначном соответствии, но с интенсиональной точки зрения это разные объекты: у них разные определения и, потенциально, разные свойства в разных контекстах. Таким образом, когда мы говорим о равной мощности, мы на самом деле имеем дело с классом эквивалентности множеств, а не с единственным объектом. Утверждение, что такие множества с совпадающими элементами обладающие одной мощностью равны, не следует автоматически - его нужно либо доказывать, либо принимать как дополнительную аксиому. Аналогичная ситуация возникает в диагональном аргументе Кантора. Диагональная последовательность $0.01(1)$ отличается от всех вещественных чисел в их позиционном представлении. Однако вещественные числа - это тоже классы эквивалентности (разные последовательности могут определять одно и то же число). Поэтому, чтобы получить противоречие, недостаточно показать, что диагональная последовательность отсутствует в списке; нужно доказать, что в списке нет ни одной последовательности из того же класса эквивалентности. Корректный подход требует явно показать, что отброшенная структурная информация (в частности, о скорости построения последовательности или о асимптотике роста длины последовательности) не влияет на результат в данном контексте, для чего необходимо явно специфицировать то, что именно было отброшено. То же самое справедливо и для понятия мощности: определяя класс равномощных множеств, мы обязаны указать, какие конкретно различия между представителями класса (интенсиональные особенности, способ построения и т.д.) несущественны для решаемых задач. В современных аксиоматических теориях множеств эта тонкость часто скрывается за неявными предположениями. Благодаря аксиоме бесконечности и убеждению, что бесконечное множество состоит исключительно из конечных элементов, слабая интерпретация биекции (существование структурного соответствия) автоматически превращается в сильную (количественное тождество множеств). Но нюанс состоит в том, что Кантор строит две независимых арифметики со своими операциями - кардинальную и порядковую, в рамках которых определяет два вида соответствия: абстрагированное от всякой структурной информации и сохраняющее вполне упорядочение. Иными словами, он корректно замыкает две алгебры типов эквивалентности и подобия, каждую со своим классом функциональных отношений и жесткой спецификацией уровня изоляции от структурной информации. Но в современных теориях мы находим лишь биекцию, которая допускает композицию с объектом задающим на ней порядок. Вспомогательный объект задаёт порядок на множестве и в случае построений Кантора, но Кантор не требует чтобы множество и такой объект принадлежали одной алгебре. Как мы увидим далее, неявные требования к уровню изоляции современных теорий, вообще говоря не соответствуют явным спецификациям Кантора, что и порождает ряд известных артефактов.

Чтобы диагонализация работала, мы должны гарантировать одновременную индексацию по позициям бесконечного количества элементов, то есть нам требуется некоторая унификация индексного множества. Кантор в своей работе добивается этого зафиксировав систему координат, что в принципе допускает обобщение на структурно ограниченные классы эквивалентности, даже в рамках слабой интерпретации взаимно однозначного соответствия. В этом отношении работа Кантора выделяется выдающейся осторожность и точностью. Если же мы не можем задать структуру и ограничить бесконечное индексное множество, диагональное построение становится невозможным. Таким образом, предположения о структуре бесконечности первичны по отношению к диагональному аргументу, они делают его обоснованным, а не обнаруживаются благодаря ему. Поскольку мы рассматриваем доказательство вне аксиоматических теорий множеств (которых ещё не существовало на момент его публикации), нужно чётко разделить два случая:

\begin{itemize}
    \item Сильная интерпретация: мы принимаем предположение, что взаимно однозначное соответствие с равными элементами означает равенство самих множеств. В этом случае множества индексов строк, столбцов и позиций диагонального элемента действительно совпадают. Диагональный элемент является вещественным числом и одновременно отличается от всех элементов списка - возникает противоречие. Доказательство работает.
    \item Слабая интерпретация: взаимно однозначное соответствие определяет лишь принадлежность к одному классу эквивалентности (одному типу). Тогда мы можем утверждать равенство только типов множеств и соответствующих друг другу элементов, но не равенство самих множеств. В этом случае нет оснований считать, что множество позиций диагонального элемента равно одновременно и множеству индексов строк, и множеству индексов столбцов. В общем случае, диагональное построение не приводит к противоречию. Корректность доказательства достигается только наличием структурных ограничений на классы эквивалентности, позволяющих исключить наличие в множестве элементов эквивалентных диагональному. 
\end{itemize}

Примечательно, что при сильной интерпретации становится естественным факт счётности всех диагональных последовательностей на множестве $C[0, 1]$, равно как и факт счётности множества предельных элементов, которые можно сопоставить таким последовательностям. Значит концепция потенциальной бесконечности несовместима с подобными представлениями и мы должны мыслить множество натуральных чисел как актуально бесконечное множество в контексте современных теорий множеств. Итак, доказательство Кантора становится корректным только в том случае, если индуктивную процедуру взаимно однозначного соответствия (с поэлементным равенством) мы интерпретируем в сильном смысле или задаёт строгие ограничения на структуру классов эквивалентности.

Хотя Кантору приписывают сильную интерпретацию тождественности множеств при биекции, многие признаки говорят об обратном. В отличие от Дедекинда, Кантор не использовал функциональную нотацию типа $\varphi(M) = N$, $\bar{\varphi}(N) = M$, $\bar{\varphi}\varphi = id$ \cite{cantor-equality, dedekind-equality}. Он не утверждал, что два множества одной мощности, состоящие из одинаковых элементов, тождественны. Индексы нужны Кантору не для того, чтобы утверждать будто множество индексов колонок тождественно множеству индексов строк, а для того, чтобы утверждать что для маркировки диагональных элементов не существует конечных индексов среди индексов координат. Тем самым Кантор подчёркивал противопоставление потенциальной и актуальной бесконечности, унаследованное из его первого доказательства. Сравнение вещественных чисел по всем позициям записи имеет смысл только благодаря понятию предела или сходимости. Без сходимости поразрядное соответствие перестаёт гарантировать тождество. Например, рассмотрим число $x = 0.1(1)$ в десятичной позиционной записи, факт сходимости позволяет нам проделать следующий трюк: $10x - 1 = x$, откуда имеем $x = 1/9$. Если же отбросить требование сходимости и оставить только взаимно однозначное соответствие цифр, то можно определить число $y = (1)1$ состоящее из бесконечного количества единиц, тогда $10y + 1 = y$ даёт $y = -1/9$, что в определённой степени парадоксально. Стоит отметить, что современная теория множеств позволяет работать с подобными конструкциями, которые называются p-адическими числами, те из них, которые имеют периодическую запись действительно ассоциируются с отрицательными рациональными числами на отрезке $[-1, 0]$ \cite{koblitz-p-adic}. Однако, данный факт лишь иллюстрирует гибкость нашего языка и способность именовать числами объекты, которые демонстрируют структурную консистентность, но не онтологическую обоснованность тождества на основании взаимно однозначного соответствия без вспомогательных ограничений. Поэтому маловероятно что Кантор, на глазах которого понятие предела приводит к возникновению позиционного представления вещественных чисел, благодаря чему поразрядное сопоставление обретает смысл сравнения, мог действительно полагать, что структурное соответствие означает тождество, а не подобие. Современная теория множеств претендует на достаточно строгую формализацию, допускающую подобные утверждения, однако, как намериваются показать авторы, такого рода обоснования строятся на неформализованных методологических допущениях. В рамках настоящей статьи, отметим, что пока интерпретация обеспечивающая равенство множеств в следствии индуктивного сопоставления элементов не является частью обоснования, доказательства диагональным построением не могут считаться корректными в общезначном смысле.

Из приведённых выше рассуждений, становится ясно, что диагональный метод Кантора, содержит ряд существенных ограничений, свойственных самому методу и требующих дополнительных предположений, а не множествам к которым его пытаются применить, по крайней мере если мы применяем его к множествам таких хорошо изученных и достоверных объектов как числа. В этом отношении он содержит некоторую неопределённость связанную с интерпретацией бесконечности, что сближает его с известными апориям Зенона. В частности данный метод не может быть использован для потенциально бесконечных числовых множеств и актуально бесконечных числовых множеств, вне контекста предположений о специальной структуре бесконечности. Корректность метода оказывается зависимой от свойств универсальности и тотальности алгебраической структуры определяющей функциональное отношение взаимно однозначного соответствия. В работе Кантора данная проблема решается введением системы координат, которую мы можем описать как множество в рамках современных представлений. Благодаря этому вывод Кантора не подлежит сомнению, если мы можем зафиксировать и ограничить структуру множества, то множество всех подмножеств не может быть структурировано подобным образом. Что действительно можно считать, доказательством несчётности множества вещественных чисел, при соответствующем определении понятия мощности. Однако, Кантор в своих зрелых работах даёт несколько иное определение понятия мощности, полностью абстрагированное от структурной информации, разделяя взаимно однозначное соответствия на два вида, принадлежащих каждый своей алгебре, что не находит отражение в последующих аксиоматиках. При этом недостаточная исследованность границ диагонального метода Кантора не означает, что все результаты полученные с его помощью могут оказаться не верными, можно лишь утверждать, что область применимости диагонального метода определяется некоторыми свойствами множеств, которыми обладают конечные множества, но, по всей видимости, обладают не все бесконечные множества и их структурные аннотации. В частности, мы пока не можем подвергнуть сомнению утверждение о несчётности вещественных чисел, ведь мы избегали рассмотрения данного доказательства с позиций современных аксиоматических теории множеств, где есть утверждения, которые можно интерпретировать как обоснование принципа взаимно однозначного соответствия в сильном смысле. Мы лишь показали, что картина классов эквивалентности Кантора отличается от картины экстенсиональной множественности современных теорий, которая при корректности и прозрачности всех допущений способна давать общезначный результат. Кроме того, есть и другие доказательства их несчётности, например, основанные на непрерывности поля вещественных чисел (лемме о вложенных отрезках), судя по всему именно они, а также авторитет Кантора служили поддержкой диагональному методу, несмотря на активную позицию многих математиков о неопределённости границ применимости данного подхода. Потрясающая вещь, которая открывается нам при анализе оригинальной версии теоремы Кантора, - невероятно близкое соседство концепций ставших современным мейнстримом и альтернативных маргинализованных идей. В классике не отягощённой современным формализмом развилка видна отчётливо и речь идёт о двух равноправных математических взглядах, где бесконечность бесструктурна и диагональное построение невозможно и где предполагается возможность иерархии бесконечного, тогда диагональный метод применим и подтверждает, что существуют достоверно различные формы бесконечного. Закономерно возникает вопрос: каким же образом могло получится так, что современная математика утратила половину описания действительности и рьяно защищает оставшуюся, не гнушаясь цензуры?

\printbibliography

 \end{document}